\documentclass[11pt,a4paper]{article}
\usepackage[T1]{fontenc}
\usepackage{lmodern}
\usepackage[margin=2.1cm]{geometry}
\usepackage{microtype}
\usepackage{amsmath,amssymb}
\usepackage{graphicx}
\usepackage{booktabs}
\usepackage{tabularx}
\usepackage{array}
\usepackage{enumitem}
\usepackage[font=small,labelfont=bf]{caption}
\usepackage{xcolor}
\usepackage{float}
\usepackage[colorlinks=true,linkcolor=blue!50!black,urlcolor=blue!50!black]{hyperref}
\setlist{itemsep=1pt,topsep=3pt}
\setlength{\parskip}{3pt}
\newcommand{\NC}{N_{\rm C}}
\newcommand{\NH}{N_{\rm H}}
\newcommand{\MH}{M_{\rm H}}
\newcommand{\MC}{M_{\rm C}}
\newcommand{\Eg}{E_{\rm g}}
\newcommand{\ok}{\textcolor{green!45!black}{done}}
\newcommand{\wip}{\textcolor{orange!80!black}{partial}}
\newcommand{\todo}{\textcolor{red!70!black}{to do}}
\newcolumntype{L}{>{\raggedright\arraybackslash}X}

\title{The grain-size-distribution (GSD) extension of CALIMA\\[3pt]
\large Where we are, and where we are going}
\author{Status report, 17 November 2026\\ \small revised after your review; EA is the default solver from 19 November;\\
  \small checks rerun on \texttt{main} on 22 November}
\date{}

\begin{document}
\maketitle
\vspace{-1.5em}
\begin{center}
\small Code: pyCALIMA \texttt{main} (GSD branch merged as PR \#20 on 19 November; checks rerun on \texttt{main}
89ee18a, with PRs \#19 and \#22, on 22 November).\\ Plan:
{\footnotesize\texttt{\textasciitilde/.claude/plans/okay-create-again-the-flickering-summit.md}}.
\end{center}

\tableofcontents

%======================================================================================================================
\section{Summary}

\paragraph{The goal.} RAMSES-CALIMA should evolve a full grain-size distribution, per composition, with the
microphysics it needs (growth, destruction, collisions, charge, hydrogenation), at a memory cost the simulations can
afford. The carbon side is the hard part: it runs from PAH molecules to micron grains, and its material changes with
size and with hydrogen content. We are building it as a \emph{minimal model}. A carrier, state or process stays only if
it acts on its own timescale or controls an observable that the next simpler model misses. Rungs M0 to M3 are each
adopted only after passing tests against benchmarks and observations. Everything is prototyped in pyCALIMA first
(Phase A), and the RAMSES implementation (Phase B) follows.

\paragraph{What changed after your review.}
\begin{itemize}
\item \textbf{Graphite and $\NC\approx400$} (Sect.~\ref{sec:materials}): every claim now has verified references, and
  two claims of the first version were wrong (Hensley \& Draine 2023 kept graphite; large-PAH spectra do not
  converge). The $\NC$ limit is argued from where the bands come from, with the counterpoints.
\item \textbf{The in-bin state $s$} is defined and derived (Sect.~\ref{sec:s}). For the other grain physics, bin-centre
  tables are off by up to a factor 2.8 for a single bin and by 42--50\,\% for evolved 8-bin distributions; the exact
  in-bin integral brings that to 4--6\,\%, one charge solve at the area node to 6--9\,\%
  (Sects.~\ref{sec:inbin}, \ref{sec:rateconv}).
\item \textbf{Grains tracked, PAHs tabulated} (Sect.~\ref{sec:hwhy}): what Montillaud, Andrews, Murga and Jones
  actually say, and the grains rescored with the PAH metric. Grain heating is as safe in equilibrium as PAH heating; the
  far-IR, visible and 3.4\,\textmu m opacity are not.
\item \textbf{A bottom-up channel} for the molecular bins (Sect.~\ref{sec:growth}): options and a recommendation, a
  calibrated source plus reaction-limited growth.
\item \textbf{The solver} (Sect.~\ref{sec:pl2m}): a full derivation, its assumptions, and a comparison with AREPO,
  FIRE and GADGET4. McKinnon et al.'s tests are rerun (Sect.~\ref{sec:mck}): EA and PL2M pass all and beat their
  piecewise-linear scheme where both run. RAMSES-MCR's moment-specific coefficients fix the frozen in-bin state under
  size-dependent drift (Sect.~\ref{sec:diallo}).
\item \textbf{A new default solver, EA} (Sects.~\ref{sec:eaf}, \ref{sec:variants}): after a survey and a benchmark
  study of variants at equal memory, the edge-aware power-law closure EA is the default of \texttt{pycalima.gsd} (same
  two scalars, about 22\,\% more cost in growth). It cuts McKinnon's growth error at 8 bins from 14 to 1.6\,\% and the
  worst A0 error from 12.1 to 7.1\,\% (4.4\,\% before the collision solver's fix of 20 November,
  Sect.~\ref{sec:binconv}), and represents distributions that end inside a bin; Gauss nodes are the default
  rate nodes. EAF, PL2M, affine and exponential closures and kernel averaging are selectable; every figure is redrawn
  with EA and PL2M (and EAF where it differs).
\item \textbf{Tables at a minimum} (Sect.~\ref{sec:nodetab}): per-radius node tables, integrated over each cell's own
  reconstruction, serve EA (or any closure) with about as many stored values as PL2M's $s$ tables.
\item \textbf{RT cross-sections in RAMSES's cool step} (Sect.~\ref{sec:rtgroups}): five fixed-$s$ pseudo-bins per bin,
  polled like ordinary bins and blended linearly in the cell's $s$. They cost 0.19\,\textmu s per cell-call, against
  4.0\,\textmu s for the exact EA average, and are within 4--6\,\% of it on the A0 group totals.
\item \textbf{Convergence of the thermal and chemical rates} with the bin number (Sect.~\ref{sec:rateconv}): a first
  look, and the full test planned for A5.
\end{itemize}

\paragraph{Rerun on 22 November.} The GSD checks were rerun on \texttt{main} after two merges that move them: the
charge solve ported from RAMSES (PR \#19: PCHIP kernels, the integer-lattice solve, LG5 as \texttt{gsd.charging}'s
default, \texttt{wide\_moments}) and a fix of the collision solver (PR \#22: the number update counted a grain that stays
in its bin twice, which cut the scheme to first order with a bias the error estimate could not see).
\begin{itemize}
\item \textbf{A0:} accretion and sputtering are bit-identical; SN dust and coagulation move by at most 0.3 points. The
  shattering case changes: at 8 bins EA goes from 4.4 to 7.1\,\%, PL2M from 7.8 to 5.7\,\%. The coarse runs moved (the
  reference only by 0.2\,\%), and they are converged in time (errmax $10^{-5}$ gives the same to 0.01\,\%). The old
  4.4\,\% was partly a cancellation: shattering now converges steadily with the bin number (EA 2.1\,\% at 16 bins and
  0.6\,\% at 32, from 3.0 and 1.7\,\%). EA's worst at 8 bins is 7.1\,\% (PL2M 12.1\,\%), above the 5\,\% line in that
  one case (Sect.~\ref{sec:binconv}).
\item \textbf{Rates at 8 bins} rise by about a point, all from shattering: the area node is within 8.8\,\% (heating)
  and 7.3\,\% (cooling), from 7.9 and 6.6\,\%; the other four cases stay at 0.7--2.0\,\% (Table~\ref{tab:areanode}).
  For a single bin the new charge solve changes mainly the large carbon bins (53\,nm--1\,\textmu m) in the PDR cell,
  whose charge distributions are wide: the Gauss nodes' worst falls from 4.2 to 1.6\,\%, the area node's 95th
  percentile from 7.8 to 7.0\,\% (Sect.~\ref{sec:inbin}).
\item \textbf{McKinnon's coagulation test} at 32 bins: the worst bin, just above the sticking threshold, is 24\,\% off,
  not 16\,\%. The old solver gives 23\,\% too once its step is cut tenfold, so the 16\,\% was a step-size artefact.
\item \textbf{Unchanged:} the PAH/grain seam (the grain side moves by $\le0.3$\,\%), \texttt{ac} against DustEM, the
  shattering test of McKinnon et al., and the RT pseudo-bins (IR means 1.3 $\to$ 1.4\,\%).
\item \textbf{Not rerun:} the variant study of Sect.~\ref{sec:variants}; its bench predates the \texttt{GSDMethod} API.
\end{itemize}

\paragraph{What is decided.}
\begin{itemize}
\item \textbf{Scheme:} two moments per bin, $N$ and $M$, with the edge-aware power-law closure EA inside each bin
  (default since 19 November; PL2M kept for comparison), and 8 bins per composition.
\item \textbf{Carbon:} one carbon population per bin on a single graded mass grid from $\NC = 20$ to 1\,\textmu m:
  \begin{itemize}
  \item 4 molecular bins below $\NC = 400$;
  \item covalent a-C nano-grains up to 20\,nm;
  \item small grains, 20--53\,nm;
  \item large CM20-like grains up to 1\,\textmu m.
  \end{itemize}
  Its material is THEMIS-like amorphous carbon, a-C(:H), at $\rho = 1.6$\,g\,cm$^{-3}$ (confirmed 22 November). Graphite
  is kept as the control, run at every rung.
\item \textbf{Optics:} per-gram values at a few radii per bin (node tables), integrated over each cell's
  reconstruction, kept per wavelength, so that the number and edges of the RT groups stay free; with PL2M the
  equivalent $s$ tables, indexed by the in-bin state $s$ (Eq.~\ref{eq:s}). The RT cross-sections in RAMSES's cool step
  use five PL2M pseudo-bins per bin (Sect.~\ref{sec:rtgroups}).
\item \textbf{Grain rates in RAMSES} (Sect.~\ref{sec:inbin}): one charge solve per bin, at its area node
  $a_* = \int a^3dN/\int a^2dN$, as many solves as today, within 8.8\,\% at 8 bins against 42--50\,\% at the bin
  centre; surface processes exact from the closure's moments; size-dependent sinks at two nodes, per moment.
\item \textbf{Top edge} (Sect.~\ref{sec:mck}): one overflow bin per composition above 1\,\textmu m, carrying $N$ and
  $M$, computed only when it holds more than $10^{-6}$ of the cell's dust mass, with a warning when it fills up; not
  a rebinned top bin (19--22 times too much mass in McKinnon's test) and not a hard stop.
\item \textbf{Hydrogen:}
  \begin{itemize}
  \item covalent grains below 53\,nm take their band gap from the local equilibrium of UV loss and
    re-hydrogenation, nothing tracked by default (22 November); a tracked $\MH$ (rung M1) is built in, compile-time
    for 0--4 bins, and switched on only if the acceptance tests (A5) need it;
  \item the molecular bins' hydrogenation is \emph{tabulated} from the equilibrium chain, not tracked (15 November):
    the state lags near transitions, but the observables RAMSES uses do not;
  \item these tables are judged on observables: carbon loss, heating, cooling, H$_2$ and the 3.4/3.3\,\textmu m proxy.
  \end{itemize}
\item \textbf{Growth of molecules:} still off by default, as in PDR codes; a bottom-up channel is now proposed
  (Sect.~\ref{sec:growth}) and needs your decision.
\end{itemize}

\paragraph{Where we are.} Table~\ref{tab:status} lists the work packages. In short:
\begin{itemize}
\item The scheme, optics, charging and heat capacity for all carbon carriers exist and are tested.
\item The hydrogen question has been settled for both grains and PAHs.
\item The first production tables for the molecular bins exist.
\item Not started yet: the RAMSES-mode operator of the GSD (A3c), the advection prototype (A4), the acceptance runs
  against observations (A5), the exporter (A6), and everything on the RAMSES side.
\end{itemize}

\begin{table}[H]
\centering\small
\caption{Work packages and their status.}\label{tab:status}
\begin{tabularx}{\textwidth}{@{}l l L l@{}}
\toprule
& work package & key result & status \\
\midrule
A0 & bin convergence of the scheme & EA with 8 bins per composition: worst error 7.1\,\% in five test cases (PL2M
      12.1\,\%; rerun 22 November); McKinnon+18's solver tests passed & \ok \\
A0b & rates over the in-bin profile & bin centre 42--50\,\% off at 8 bins; EA exact 4.4--6.2\,\%, Gauss nodes $\le$\,7.1\,\%;
      RAMSES: one charge solve per bin at the area node, $\le$\,8.8\,\% (decided 18 November) & \ok \\
A0c & selectable solver methods & \texttt{pycalima.gsd.closures}: EA default, EAF/PL2M/affine/exp, Gauss nodes, kernel
      averaging; node tables for any closure & \ok \\
A1 & THEMIS CM20 optics & reproduces DustEM's J13 extinction to 0.9\,\%; Henke extension below 40\,nm & \ok \\
A2 & per-bin optics tables $\kappa_i(\lambda,s)$ & every observable within 0.5\,\% & \ok \\
A3a & materials, graded grid, $\rho = 1.6$ & material registry (unknown names raise), \texttt{CarbonLayout} & \ok \\
A3b & per-carrier tables & a-C(:H) optics on an $\Eg$ axis; charging \texttt{ac}/\texttt{ach}/\texttt{cm20}, 5 intermediate nodes;
      heat capacity; PAH/grain seam & \wip \\
step 1 & tracked-hydrogen machinery & surface H, size-limited gap ($\le$\,1\,\% heating with 5 nodes), H-axis tables & \ok \\
step 2 & tabulate or track (grains) & E in lit gas, M in transition layers and shielded gas & \ok \\
step 2 & tabulate or track (PAHs) & tables, judged on observables; production tables built (deep chain, +8) & \ok \\
A3c & \texttt{ramses} mode of \texttt{gsd} & the per-cell operator with the $\MH$ block and the exact H step & \todo \\
A4 & advection prototype & ratio pass, donor-ratio TVA flux, $\MH$ as a child of $\MC$ & \todo \\
A5 & accept or reject & DustEM SED/EXT, diffuse ISM, PDR one-zone and slab runs; ladder ablations & \todo \\
A6 & exporter & all per-bin tables for RAMSES with provenance & \todo \\
B & RAMSES & state layout, sources, transport, tables (ask before branching) & \todo \\
\bottomrule
\end{tabularx}
\end{table}

\paragraph{The next steps,} in order (Sect.~\ref{sec:next}):
\begin{enumerate}
\item consolidate: done for the GSD branch (merged as PR \#20; checks rerun on \texttt{main} on 22 November); merge
  pahchem's PRs \#17 and \#18 and rerun the checks that read pahchem;
\item settle the open choices: the bottom-up channel, the in-bin treatment of the rates, the seam, the super-H
  physical stop, the $G_0$ axis above $10^4$;
\item build A3c, the GSD operator RAMSES will run per cell;
\item A4 and A5, the advection prototype and the acceptance tests;
\item A6, the exporter, then Phase B in RAMSES.
\end{enumerate}

%======================================================================================================================
\section{The design in one picture}

Figure~\ref{fig:network} shows the carriers and what they exchange. The rules behind it:
\begin{itemize}
\item \textbf{One carbon grid.} Every carrier is carbon on one mass grid, with $\NC = m/m_{\rm C}$ exact. The radius
  follows from $\rho = 1.6$\,g\,cm$^{-3}$ wherever geometry enters. A bin's size and hydrogen content decide what it is:
  a PAH, a hydrogenated nano-grain or a-C.
\item \textbf{What is stored.} Each bin stores $N$ and $\MC$. Covalent bins below 53\,nm can also store $\MH$ (M1),
  advected as a child of $\MC$. Charge, dust temperature, band emission and the molecular-bin hydrogen state are tables.
\item \textbf{Moving carbon moves hydrogen.} Every process that moves carbon between bins also moves hydrogen at the
  source bin's H/C.
\item \textbf{Conditions are inputs.} $G_0$, $n$(H), $T$ and $n_e$ come from the simulation or from sourced environment
  cards, never from a model inside the dust code. Anything tuned stays opt-in and labelled ad hoc.
\end{itemize}

\begin{figure}[H]
\centering
\includegraphics[width=\textwidth]{figs/carbon_network.png}
\caption{The carbon carriers of the plan and their exchanges with each other, the radiation and the gas. The badges
mark the rung that adds each process (M0 is the baseline); dashed boxes are optional rungs. \emph{Out of date since
23 November:} the PAH-cluster box (M2) was dropped from the model and photo-fragmentation is an opt-in process
(Sect.~\ref{sec:transfer}); the figure's source is updated, the image is to be re-rendered.}
\label{fig:network}
\end{figure}

\begin{table}[H]
\centering\small
\caption{The minimal-model ladder. Each rung is run with and without its addition, and adopted only if it brings an
observable inside its observed range where the rung below fails. Memory is for $10^6$ cells, with the TVA drift stored or
recomputed (B4).}\label{tab:ladder}
\begin{tabularx}{\textwidth}{@{}l L l l@{}}
\toprule
rung & adds & carbon scalars & memory per $10^6$ cells \\
\midrule
M0 & the GSD with THEMIS-like carbon; molecular tables below $\NC = 400$; equilibrium H everywhere & 22 & 620 / 350 MB \\
M1 & tracked $\MH$ in the covalent bins below 53\,nm: UV H loss, re-hydrogenation, H$_2$ release, $\Eg$ follows H/C &
  +4 (was +8) & $\approx$ 680 / 420 MB \\
M2 & \emph{removed 23 November} (no PAH clusters); a-C(:H) photo-fragmentation is now an opt-in process & -- & \\
M3 & a hydrogenated skin on large grains (dense-gas mantles) & +3 & \\
\bottomrule
\end{tabularx}
\end{table}
\noindent Since the molecular bins are now tabulated, M1 needs $\MH$ only in the 4 covalent bins below 53\,nm, not in
8 bins as planned. Its memory is interpolated from the plan's 740\,/\,480\,MB for the 8-bin version.

%======================================================================================================================
\section{The size-distribution solver: two moments and edge-aware power laws}\label{sec:pl2m}

The solver stores two moments per bin, $N$ and $M$, and rebuilds the distribution inside each bin from them. Its
default method since 19 November is EA: a power law in each bin (PL2M) except in bins at an edge of the distribution,
which are treated as fronts (Sect.~\ref{sec:eaf}). PL2M stays available as a method for comparison, together with
EAF and the affine and exponential closures (\texttt{pycalima.gsd.closures}). This section derives the scheme step by
step: what is evolved, how a bin is represented, how each process acts on
that representation, how the bins are transported in RAMSES, which assumptions it rests on, and how it differs from the
GSD solvers of AREPO, GIZMO/FIRE and GADGET. The code is \texttt{pycalima.gsd}. The design review behind it, with the
numerical literature (aerosol and planet-formation schemes), is \texttt{docs/development/gsd\_solver\_review.md}; the
comparison with galaxy codes is in \texttt{results/gsd\_overview/lit/gsd\_solvers.md}.

\subsection{What is evolved}
For each composition (carbon, silicate) let $f(m,\mathbf x,t)$ be the number of grains per unit volume and per unit
grain mass. It obeys
\begin{equation}\label{eq:gde}
\frac{\partial f}{\partial t} + \nabla\!\cdot\!\big[f\,(\mathbf v + \mathbf w(m))\big]
 + \frac{\partial (\dot m f)}{\partial m} = \mathcal C[f] + S(m),
\end{equation}
where $\mathbf v$ is the gas velocity, $\mathbf w(m)$ the drift of grains of mass $m$ relative to the gas,
$\dot m = 4\pi a^2\rho\,\dot a$ the growth (accretion) or loss (sputtering, sublimation) of a single grain,
$\mathcal C[f]$ the collisions (coagulation, cratering, shattering) and $S$ the sources and sinks (stellar injection,
nucleation, a bottom-up PAH source, carbon loss of the molecules). Mass and radius are tied by
$m = \tfrac43\pi\rho a^3$ ($\rho = 1.6$\,g\,cm$^{-3}$ for carbon, 3.3 for silicate), and $\NC = m/m_{\rm C}$.

\subsection{Bins, moments and the closure problem}
The mass axis is cut into fixed bins $I_i = [m_i, m_{i+1}]$ with mass ratio $\eta_i = m_{i+1}/m_i$ (for carbon,
$\eta = 2.1$ in the four molecular bins and 18.8, a factor 2.66 in radius, in the seven grain bins). Each bin stores
two numbers per cell:
\begin{equation}
N_i = \int_{I_i} f\,dm, \qquad M_i = \int_{I_i} m f\,dm .
\end{equation}
Integrating Eq.~(\ref{eq:gde}) over $I_i$ (transport aside, and using $\int m\,\partial_m(\dot m f)\,dm = [m\dot m f] -
\int \dot m f\,dm$):
\begin{align}
\frac{dN_i}{dt} &= (\dot m f)_{m_i} - (\dot m f)_{m_{i+1}} + \int_{I_i}(\mathcal C + S)\,dm, \label{eq:dN}\\
\frac{dM_i}{dt} &= \int_{I_i}\dot m f\,dm + m_i(\dot m f)_{m_i} - m_{i+1}(\dot m f)_{m_{i+1}}
 + \int_{I_i} m\,(\mathcal C + S)\,dm. \label{eq:dM}
\end{align}
The right-hand sides need $f$ at the bin edges and inside the bins, which $(N_i, M_i)$ do not determine. A scheme is
defined by its \emph{closure}, the shape $f$ it assumes inside a bin given the stored moments, and by how each process
uses that shape. One-moment schemes store only $M_i$ and must also assume $N_i$, i.e.\ fix the shape a priori; A0
tested two such closures (an MRN slope inside every bin, and slopes from the neighbours).

\subsection{The power law inside a bin, and the in-bin state \texorpdfstring{$s$}{s}}\label{sec:s}
PL2M closes the system with a power law in each bin, $f(m) = A_i\,m^{\beta_i}$ on $I_i$. Then
\begin{equation}
N_i = A_i m_i^{\beta+1}\frac{\eta^{\beta+1}-1}{\beta+1},\qquad
M_i = A_i m_i^{\beta+2}\frac{\eta^{\beta+2}-1}{\beta+2},
\end{equation}
and the mean mass $\bar m_i = M_i/N_i$ fixes the exponent through
\begin{equation}\label{eq:R}
\frac{\bar m_i}{m_i} = R(\beta;\eta) \equiv \frac{\beta+1}{\beta+2}\,\frac{\eta^{\beta+2}-1}{\eta^{\beta+1}-1}
\end{equation}
(with the limits $R = (\eta-1)/\ln\eta$ at $\beta=-1$ and $R = \eta\ln\eta/(\eta-1)$ at $\beta=-2$). $R$ rises
monotonically from 1 ($\beta\to-\infty$, every grain at $m_i$) to $\eta$ ($\beta\to+\infty$, every grain at $m_{i+1}$),
so every \emph{realisable} pair, $m_i \le M_i/N_i \le m_{i+1}$, has exactly one $\beta$. It is found by a bracketed
(Illinois) iteration in about 40\,\textmu s; $A_i$ then follows from $N_i$. In radius the profile is $dN/da \propto
a^{\lambda}$ with $\lambda = 3\beta+2$. Bins steeper than $|\beta| = \min(150, 300/\ln\eta)$ are held as spikes,
$N_i$ grains at $\bar m_i$.

\paragraph{The in-bin state.} We define
\begin{equation}\label{eq:s}
s_i = \frac{\bar m_i/m_i - 1}{\eta_i - 1} \in [0,1], \qquad \bar m_i = M_i/N_i .
\end{equation}
$s = 0$ means every grain sits at the lower edge, $s = 1$ every grain at the upper edge. Since $R$ depends on $\beta$
and $\eta$ only, $s$ fixes the whole shape inside the bin (Fig.~\ref{fig:pl2m}a,b). Any per-gram average of a
single-grain property $k(a)$ (an opacity, a heating rate per grain over its mass, ...),
\begin{equation}\label{eq:kavg}
\langle k\rangle_i(s) = \frac{\int_{I_i} k(a)\,m\,dN}{\int_{I_i} m\,dN},
\end{equation}
therefore depends on the cell only through $s$ and the local conditions. That is what the optics tables
$\kappa_i(\lambda, s)$ use (Sect.~\ref{sec:materials}). $s$ describes the population inside a bin, not a size: two cells
with the same $s_i$ have the same in-bin shape and the same per-gram rates. Reference values: an MRN slope
($dN/da\propto a^{-3.5}$, $\beta = -11/6$) gives $s = 0.14$ in a grain bin and 0.39 in a molecular bin; equal number
per $\ln m$ ($\beta=-1$) gives 0.29 and 0.44; equal mass per $\ln m$ ($\beta=-2$) gives 0.12 and 0.38.

\paragraph{Why a power law.}
\begin{itemize}
\item It is positive for every realisable $(N, M)$, so it needs no slope limiter.
\item It is local: it needs no neighbouring bin, unlike slope-limited linear reconstructions.
\item It covers the whole realisable range, down to a population at a single size, so a nucleation burst or a
  narrow injected population stays sharp.
\item Integrals of powers of $a$ (moments, power-law fragment spectra) are analytic.
\item A scale-free distribution such as MRN is exact inside every bin.
\end{itemize}
The price: a bin cannot hold structure narrower than itself (two peaks, a cut-off inside the bin), and the
reconstruction jumps at bin edges (Fig.~\ref{fig:pl2m}c).

\subsection{The default closure: edge-aware power laws (EA)}\label{sec:eaf}
A power law fills the whole bin. Where the distribution ends inside a bin, a growth front or the edge of a narrow
population, PL2M must still spread the bin's $N$ and $M$ over the bin and piles grains against one edge
(Sect.~\ref{sec:mck}); after every step this wrong shape is moved and re-imposed, which is PL2M's numerical diffusion.
The default closure since 19 November, EA, keeps the two stored moments and changes only the shape in bins that sit at
an edge of the distribution (Sect.~\ref{sec:variants} explains how it was chosen):
\begin{itemize}
\item \textbf{Is a neighbour empty?} For bin $i$ and each neighbour, $q = (N_{\rm nb}/\Delta a_{\rm nb})/(N_i/\Delta a_i)$
  compares the number per unit radius; the ends of the grid count as empty. A neighbour is a front's empty side to the
  degree $e = \max(0, 1 - q/\epsilon)$, with $\epsilon = 10^{-6}$, so smooth tails are never mistaken for fronts.
\item \textbf{The profile} is a blend of four shapes, each carrying the bin's own $N_i$ and $M_i$:
  \begin{equation}
  f_i = (1-e_L)(1-e_R)\,f_{\rm PL2M} + e_R(1-e_L)\,f_{R} + e_L(1-e_R)\,f_{L} + e_Le_R\,N_i\,\delta(m - \bar m_i).
  \end{equation}
  With the right neighbour empty, $f_R\propto a^{\lambda_f}$ on $[a_i, a_c]$: the distribution continues into the bin
  with the slope of its other (left) neighbour, $\lambda_f = \lambda_{i-1}$, and ends at the cut $a_c$ that makes its
  mean mass $\bar m_i$; $f_L$ is the mirror image, on $[a_c, a_{i+1}]$ with $\lambda_f = \lambda_{i+1}$ (flat,
  $\lambda_f = 0$, if that neighbour is itself a front). With both neighbours empty the bin is held at its mean mass
  (a moving pivot). Where no cut fits, the PL2M power law is used.
\item \textbf{A power-law segment with sharp ends is exact}: the fronts continue the segment's own slope, so a boxcar,
  an MRN with a cut-off or a narrow population are represented as they are, and growth moves them without diffusion.
\item \textbf{The variant EAF} fixes the front exponent at $\lambda_f = -3.5$ (an MRN slope), so that every shape
  depends on $s$ alone and one-dimensional $s$ tables suffice (four per bin). It was the default for a day: on flat
  profiles its fixed slope puts peaks at the edges (Figs.~\ref{fig:mck9}, \ref{fig:mckgrowth}). EA's tables are kept small
  instead with per-radius node tables (Sect.~\ref{sec:nodetab}).
\item \textbf{What does not change:} $N$ and $M$ are exact, the profile is positive and realisable, and the remap,
  the collision nodes and the in-bin integrals work on these shapes exactly as on a power law. The drift velocities use
  the local power law, never the pivot, which would freeze the in-bin segregation.
\end{itemize}
PL2M stays available as a method (\texttt{closure="pl2m"}), together with EAF and the affine and exponential closures,
for comparison.

\begin{figure}[H]
\centering
\includegraphics[width=\textwidth]{figs/pl2m_scheme.pdf}
\caption{The scheme. (a) Inside a grain bin the profile is a power law fixed by $s$; the ticks mark the two
nodes the collision operator uses (Sect.~\ref{sec:coll}). (b) $s$ against the exponent $\beta$ for the molecular and
grain bins, with the MRN slope and the flat-in-number and flat-in-mass profiles marked. (c) A real case: silicate
after 5\,Myr of shattering, coagulation and accretion of SN dust in the WNM (A0), EA and PL2M with 8 bins against
the reference (identical here: no bin is at an edge; dotted lines: the bin edges). (d) One remap step: an $a^{-3.5}$
distribution that ends inside a bin, grown by a constant $\Delta a$. The exact image (dashed) is cut at the fixed edges
and stored as $(N_i, M_i)$; EA rebuilds both ends as fronts where they are, PL2M spreads them over the end bins.}
\label{fig:pl2m}
\end{figure}

\subsection{Growth and erosion along characteristics}
Rather than evaluating the edge fluxes $\dot m f$ of Eqs.~(\ref{eq:dN})--(\ref{eq:dM}) (an Eulerian upwind scheme,
which is diffusive and limited by a CFL condition set by the smallest grains; Hirashita \& Aoyama 2019 report the
smoothing it causes at 16 bins), the solver moves the grains along their characteristics. With the conditions frozen over a
step $\Delta t$, a grain of radius $a$ ends at $X(a)$, the solution of $da/dt = \dot a(a)$. Number is conserved along
characteristics, $f(X)\,dX = f(a)\,da$, so
\begin{equation}
N_j^{n+1} = \sum_i \int_{a\in I_i,\;X(a)\in I_j} \Big(\frac{dN}{da}\Big)_i da,\qquad
M_j^{n+1} = \sum_i \int_{a\in I_i,\;X(a)\in I_j} m\big(X(a)\big)\Big(\frac{dN}{da}\Big)_i da .
\end{equation}
In practice (\texttt{gsd.remap}):
\begin{itemize}
\item $\dot a(a)$ is tabulated on a fine logarithmic grid, and $X(a)$ is integrated exactly for an $\dot a$ linear
  between table points. The step has no time-step limit, and equilibrium radii ($\dot a = 0$, e.g.\ in sublimation)
  are approached but never crossed.
\item Each bin is split into four pieces in $\ln a$ whose images are linear in $a$. This is exact for a constant
  $\dot a$, which is the case for accretion of neutral species and for thermal sputtering at a fixed yield.
\item The images are cut at the fixed edges, and each piece's number and mass are integrated analytically against
  $u^\lambda$ ($u = a/a_i$; the mass integrand $X(u)^3 u^\lambda$ is a cubic times a power).
\item Grains mapped below the smallest edge are removed and their mass is returned to the gas; grains above the
  largest edge go to an overflow.
\end{itemize}
Number is moved exactly, mass is exact for the flow map, the map is monotone (characteristics do not cross), and there
is little numerical diffusion. Figure~\ref{fig:pl2m}d shows the cost: after the step the shape inside each bin is a
power law again. A spike inside a bin, re-represented as a steep power law, loses about 0.5\,\% of its mass per
factor 10 of growth.

\paragraph{Exchange with the gas and sources} (\texttt{gsd.exchange}). Accretion and sublimation exchange a key
species X with the gas. The step is implicit in the gas: the end-of-step density $n_X^{n+1}$ is the root of
\begin{equation}
n_X^{n+1} + X_{\rm dust}\big(n_X^{n+1}\big) = X_{\rm tot},
\end{equation}
where $X_{\rm dust}$ is what the remap leaves on the grains when the growth rate is evaluated at $n_X^{n+1}$. The left
side is monotone in $n_X^{n+1}$, so the root is unique and positive, and every element is conserved to round-off (the
idea of Jacobson's analytical predictor of condensation). Nucleation, where it matters (AGN rims, SN ejecta), is
injected along the characteristic from the critical size; stellar injection (SN, AGB) adds its source spectrum
integrated exactly over the bins; a bottom-up PAH source (Sect.~\ref{sec:growth}) would inject into the first
molecular bin.

\subsection{Collisions}\label{sec:coll}
Coagulation and shattering follow the Smoluchowski equation with fragmentation,
\begin{equation}
\frac{\partial f}{\partial t}\Big|_{\rm coll} = \frac12\int_0^m K(m',m-m')\,f(m')f(m-m')\,dm'
 - f(m)\int_0^\infty K(m,m')f(m')\,dm' + (\text{fragments}),
\end{equation}
with the kernel $K = \sigma\,\Delta v\,P\times$ (Coulomb factor). The solver evaluates the pair integrals with \textbf{two
nodes per bin} (\texttt{collisions.bin\_nodes}). The bin's number is split at its median,
\begin{equation}
\int_1^{u_{\rm med}} u^\lambda du = \tfrac12\int_1^{u_R}u^\lambda du \;\Rightarrow\;
u_{\rm med} = \big[\tfrac12(1+u_R^{\lambda+1})\big]^{1/(\lambda+1)},\qquad u_R = \eta^{1/3}.
\end{equation}
Node 1 carries $N_i/2$ grains at the mean mass of the lower half,
$m_{i,1} = m_i\int_1^{u_{\rm med}}u^{\lambda+3}du\big/\int_1^{u_{\rm med}}u^{\lambda}du$; node 2 carries $N_i/2$ at
$m_{i,2} = 2M_i/N_i - m_{i,1}$. The two nodes reproduce $N_i$ and $M_i$ exactly.
\begin{itemize}
\item \textbf{Pairs.} Node pairs $(p,q)$ collide at $r_{pq} = K_c(p,q)\,w_p w_q/2$ per outcome channel $c$
  (coalescence, cratering, shattering), with $w$ the node weights.
\item \textbf{Products.} The remnant (merged grain, or cratered target) is deposited as one grain spread uniformly in
  mass over its parents' node intervals; a point deposit would switch bins discontinuously as a node crosses an edge.
  Fragments with $dN/da\propto a^{-\alpha_f}$ on $[a_{\rm lo}, a_{\rm hi}]$ are integrated exactly over the bins, at
  $O(1)$ cost per pair (difference arrays; the idea of Rafikov et al.\ 2020).
\item \textbf{Velocities and charges.} Brownian, turbulent (Ormel \& Cuzzi 2007; Yan et al.\ 2004 as in Hirashita \&
  Yan 2009) and drift velocities; Coulomb factors from the charge distributions (three Gauss--Hermite charge points
  per radius).
\item \textbf{Time integration.} The mass transfers form a production--destruction system,
  $dM_k/dt = \sum_j [p_{kj}(M) - d_{kj}(M)]$ with $p_{kj} = d_{jk}$, integrated with MPRK22 (Kopecz \& Meister 2018):
  \begin{align*}
  M_k^{(1)} &= M_k^n + \Delta t\sum_j\Big[p^n_{kj}\frac{M_j^{(1)}}{M_j^n} - d^n_{kj}\frac{M_k^{(1)}}{M_k^n}\Big],\\
  M_k^{n+1} &= M_k^n + \frac{\Delta t}{2}\sum_j\Big[(p^n_{kj}+p^{(1)}_{kj})\frac{M_j^{n+1}}{M_j^{(1)}}
     - (d^n_{kj}+d^{(1)}_{kj})\frac{M_k^{n+1}}{M_k^{(1)}}\Big].
  \end{align*}
  Each stage is one linear solve with an M-matrix, so the masses stay positive for any $\Delta t$, mass is conserved to
  round-off, and the scheme is second order. Number losses are Patankar-weighted in each bin (positive), and number
  gains follow the delivered mass. The difference between the first- and second-order stages sets the sub-steps.
\item \textbf{Options.} The nodes are computed from the bin's reconstruction (EA by default); a 2-point Gauss rule
  in radius (\texttt{nodes="gauss"}) and kernel averaging over sub-slices of each node pair (\texttt{kernel\_average})
  are available (Sect.~\ref{sec:variants}).
\item \textbf{Checks.} The constant, additive and multiplicative Smoluchowski solutions, and Hirashita \& Yan (2009)
  shattering; about 2\,ms per sub-step on 150 bins.
\end{itemize}

\subsection{Transport in RAMSES (A4: designed, not built)}
$N_i$ and $M_i$ are advected as passive scalars with the gas, and each bin drifts with its terminal velocity $\mathbf
w_i$ (TVA):
\begin{equation}
\partial_t N_i + \nabla\!\cdot[N_i(\mathbf v + \mathbf w_i)] = 0,\qquad
\partial_t M_i + \nabla\!\cdot[M_i(\mathbf v + \mathbf w_i)] = 0 .
\end{equation}
The constraint is realisability, $m_i \le M_i/N_i \le m_{i+1}$. Mixing two realisable states of one bin gives
$\bar m = (N_1\bar m_1 + N_2\bar m_2)/(N_1+N_2)$, which lies between $\bar m_1$ and $\bar m_2$: the realisable states
form a convex cone. Any update that writes a cell's new $(N_i, M_i)$ as a positive combination of old states
(first-order upwind under the CFL limit, AMR restriction) keeps them realisable. A MUSCL scheme with independent slopes
for $N$ and $M$ can break this at the cell faces. The design therefore:
\begin{itemize}
\item advects $M_i$ with the gas scheme, inside the CMA group of its element, so that the groups still sum to the
  element;
\item gives $N_i$ the donor-ratio flux $F_N = F_M/\bar m_{\rm face}$, with $\bar m_{\rm face}$ reconstructed with a
  limited slope, so that it lies between the neighbouring cells' values and hence inside the bin;
\item treats $\MH$ the same way, as a child of $\MC$ with the ratio $x = \MH/\MC$ (as CMA treats ion fractions);
\item moves $N_i$ and $M_i$ with their own drift velocities, $w^N_i = \int w\,dN/N_i$ and $w^M_i = \int w\,m\,dN/M_i$,
  from the bin's power law (the moment-specific coefficients of Diallo et al.\ 2026 and Hopkins 2023;
  Sect.~\ref{sec:diallo}). The packet that crosses a face then has the mean mass $\int m\,w\,dN/\int w\,dN$, which
  lies inside the bin, so realisability holds, and big and small grains of one bin can separate ($s$ changes); with one
  velocity per bin they cannot;
\item injects the parent's ratios at AMR prolongation, and sums $N$ and $M$ at restriction;
\item keeps a realisability clip as a counted safety net, whose count A4 will report.
\end{itemize}

\subsection{One RAMSES step}
After the hydro step, \texttt{dust\_solver.f90} advances each cell with its conditions frozen:
\begin{enumerate}
\item the exact hydrogen step of the bins that carry $\MH$ (rung M1);
\item collisions, in MPRK22 sub-steps;
\item sources: stellar injection, the bottom-up source if on, carbon loss of the molecular bins from their tables;
\item the exchange with the gas (accretion, sputtering, sublimation), last because it is the stiffest coupling
  (Sportisse 2000).
\end{enumerate}
The order is to be fixed in A3c, where the per-cell operator is built and timed.

\subsection{Assumptions}
\begin{enumerate}
\item Grains are spherical, one material per population; the radius follows from $\rho$, and $\NC = m/m_{\rm C}$.
\item Inside a bin the distribution is a power law in mass; structure narrower than a bin is lost.
\item The conditions in a cell (density, temperature, radiation, electron density) are uniform and frozen over a step;
  collisions are sub-stepped inside it.
\item Grain charges are in equilibrium (they relax in hours to weeks), from the charge solve in each cell.
\item Collisions see each bin through its two nodes. Grains of different compositions do not collide (no mixed grains).
\item Grains move with the gas plus a terminal drift per bin: the stopping time is short against the cell-crossing and
  dynamical times, and grains of different sizes inside one bin share its drift.
\item The grid is fixed, with an overflow above 1\,\textmu m. Fragments below the grid go to the gas or to the smallest
  bin, a recorded choice.
\item Stochastic heating enters only the molecular tables and post-processing; the solver uses one charge distribution
  and one rate per radius.
\end{enumerate}

\subsection{How it compares with the GSD solvers of other codes}\label{sec:codes}
Every full-GSD solver in a galaxy code descends from Hirashita \& Yan (2009), through Asano et al.\ (2013) and
McKinnon et al.\ (2018). Three lines are in use today: McKinnon's in AREPO, its port to FIRE by Choban et al.\ (2026),
and the Osaka scheme of Hirashita \& Aoyama (2019) in GADGET4 (Aoyama et al.\ 2020; Romano et al.\ 2022; Matsumoto et
al.\ 2026). Table~\ref{tab:codes} compares them with our scheme feature by feature; the papers and quotes are in
\texttt{lit/gsd\_solvers.md}, and Sect.~\ref{sec:mck} reruns McKinnon's tests.

\begin{table}[H]
\centering\footnotesize
\caption{The three GSD solvers in use in galaxy codes, against ours (EA, the default). McKinnon+18: AREPO; Choban+26: FIRE (GIZMO);
GADGET4: Matsumoto+26 with the scheme of Hirashita \& Aoyama (2019), Aoyama+20 and Romano+22.}\label{tab:codes}
\begin{tabularx}{\textwidth}{@{}>{\raggedright\arraybackslash}p{2.3cm} L L L L@{}}
\toprule
 & McKinnon+18 (AREPO) & Choban+26 (FIRE) & GADGET4-OSAKA & EA (this work) \\
\midrule
stored per bin & $N$ and slope ($\equiv N$, $M$) & as McKinnon & mass only & $N$, $M$ \\
shape in a bin & linear in $a$; slope limiter keeps $M$, changes $N$ & as McKinnon & constant & power law in $m$, fronts at the edges of the distribution (EA); positive by construction \\
bins, range & 16; 1\,nm--1\,\textmu m & 16; 5\,nm--1\,\textmu m & 30; 0.3\,nm--10\,\textmu m & 8 per composition; 0.35\,nm--1\,\textmu m \\
compositions & one GSD for all dust & one per species (silicate, carbon, Fe) & one; composition in post-processing & one per composition, carbon graded (molecular bins) \\
growth, sputtering & shift by $\dot a\Delta t$ + intersection with the bins & as McKinnon; $\Delta t\le0.3\,\delta a_1/\dot a$ & first-order upwind in mass (diffusive) & characteristic remap, any $\dot a(a)$, no step limit \\
past the grid edges & rebinned into the edge bin (mass kept) & as McKinnon & -- (grid to 10\,\textmu m) & removed to the gas (bottom), overflow (top) \\
collisions & HY09 kernels on the linear profiles, explicit first order, masses rescaled & as McKinnon, sub-cycled on mass moved and number lost & Smoluchowski sums, loss semi-implicit & two nodes per bin, MPRK22 (positive and conservative for any step) \\
gas coupling & from neighbour cells, capped by available metals & in the cell & in the particle & implicit in the gas, elements to round-off \\
transport & dust particles with drag; one stopping time per GSD & with the gas cell + turbulent diffusion & with the SPH particle + turbulent diffusion & passive scalars + drift per bin and moment \\
sub-grid ISM & -- & clumping factors, ad hoc coagulation boost & two phases per particle & conditions are cell inputs \\
convergence shown & growth: 1st/2nd order; collisions vs a many-bin solver & GSD at 16 bins, extinction near 8 & HA19: 16 vs 128 bins diffusive & A0 observables at 8 bins; McKinnon's tests (Sect.~\ref{sec:mck}) \\
\bottomrule
\end{tabularx}
\end{table}

\paragraph{McKinnon against FIRE.} Choban et al.\ (2026) say their method ``largely follows the second-order
piecewise linear method developed in McKinnon et al.\ (2018)'': the same profile, slope limiter, shift-and-intersect
update and top-edge rebinning (their Eqs.~E5--E11 are McKinnon's Eqs.~34--46). What changes is where the GSD lives (in
each gas cell, not on dust particles), one GSD per species, explicit sub-cycling criteria (they report that without
them coagulation ``heavily reduced'' the large grains), and sub-grid factors: clumping, Coulomb enhancement of accretion
and an ad hoc coagulation boost.

\paragraph{McKinnon against GADGET4.} The Osaka line stores only the mass per bin with a constant profile, moves
growth and sputtering by first-order upwind, and needs many bins (30, from 0.3\,nm to 10\,\textmu m); HA19 show the
numerical diffusion at 16 bins. Its grid ends at 10\,\textmu m, so the top-edge rebinning that McKinnon's Fig.~11
warns about does not arise. Composition is not evolved: it is assigned in post-processing. Each particle carries two
sub-grid phases, and the processes are split between them.

\paragraph{Where our scheme stands.}
\begin{itemize}
\item It shares McKinnon's and FIRE's two moments and characteristic update, but its profile is a power law in $m$,
  which never goes negative, so no limiter changes $N$ (McKinnon's slope limiter keeps $M$ at the cost of $N$); and
  with EA it represents a distribution that ends inside a bin, which none of the others does.
\item Its collisions are positive and conservative for any step (MPRK22), where the others take explicit or
  semi-implicit first-order steps with mass rescaling and sub-cycling; we found no dust code that uses MPRK22.
\item Its growth step takes any size-dependent $\dot a(a)$ (sublimation, Kelvin term), not only a constant shift.
\item It uses 8 bins per composition (16 in AREPO and FIRE, 30 in GADGET4), and the in-bin state $s$ as a table axis.
\item It is the only one that moves $N$ and $M$ as passive scalars with a size-dependent drift in an Eulerian code;
  the only terminal-drift treatment in RAMSES (Lebreuilly et al.\ 2019) has fixed sizes, and McKinnon's particles
  share one stopping time per GSD.
\item RAMSES-CALIMA today: two sizes per composition (Dubois et al.\ 2024; Rodr\'iguez Montero et al.\ 2026). Its paper
  (Sect.~5.2) flags that processes are evaluated at the central size of each bin, that memory grows with the bins, and
  that large grains decouple in hot, low-density gas; the in-bin treatment (Sect.~\ref{sec:inbin}) and the drift
  address the first and last.
\end{itemize}

%======================================================================================================================
\section{McKinnon et al.'s solver tests, rerun with EA and PL2M}\label{sec:mck}

McKinnon et al.\ (2018, Sect.~3) test their scheme on a set of problems with known answers. We reran all of them
(their Figs.~9--13 and 15--16) with the default method, EA, with PL2M for comparison and with EAF where it differs,
in
\texttt{diagnostics/gsd/check\_mckinnon18\_tests.py}. For the growth tests we also re-implemented their
piecewise-constant (PC) and piecewise-linear (PL) schemes from their Sect.~3.2, so that all schemes run on the same
problem with the same steps; at 64 bins our PC and PL give 34\,\% and 1.2\,\% mass
error where they quote about 40\,\% and 1\,\%, which validates the re-implementation. The collision tests are compared
with an independent fine solution (640 bins, scipy LSODA), as theirs were. Results are in
\texttt{figs/mckinnon18\_tests.txt}.

\subsection{One step of the scheme (their Fig.~9)}
Figure~\ref{fig:mck9} redraws their schematic with the five schemes on the same few bins: a boxcar that ends inside
a bin, before and after a growth step. PL needs its slope limiter where the boxcar's edge enters a bin (the profile
would go negative); the limiter keeps the bin's mass and changes its number. PL2M needs no limiter, but in a
partly filled bin its power law piles the grains against one edge: number and mass are exact, the local shape is not.
EA puts the edge inside the bin where it is and reproduces the boxcar exactly. EAF also places the edge, but its fixed
front exponent ($-3.5$) slopes where the boxcar is flat and peaks at the cut: this is why EA, not EAF, is the default.

\begin{figure}[H]
\centering
\includegraphics[width=\textwidth]{figs/mckinnon18_fig9.pdf}
\caption{After McKinnon et al.'s Fig.~9: a boxcar (grey) on five bins at time $t$ and after all grains have grown by
$\Delta a$, as PC (dashed), PL (green), PL2M (orange), EAF (purple) and EA (blue, the default) represent it. Dotted
lines: bin edges.}
\label{fig:mck9}
\end{figure}

\subsection{Growth: convergence and a growth--shrinking cycle (their Figs.~10, 13)}
\begin{itemize}
\item \textbf{Convergence} (Fig.~\ref{fig:mckgrowth}a): a boxcar from 1\,nm to 5.6\,nm grown at 0.005\,\textmu m/Gyr
  for 5\,Gyr, in 100 steps, against the exact shifted boxcar. PC is first order and PL second order, as they found.
  EA is 1.6\,\% at 8 bins, 0.5\,\% at 16, 0.15\,\% at 32 and exact to round-off from 64; PL2M 14\,\%, 4.6\,\%, 1.2\,\%
  and 0.25\,\%; PL 51\,\%, 18\,\%, 4.8\,\% and 1.2\,\% (EAF: 2.2\,\%, 0.5\,\%, 0.11\,\%, 0.013\,\%). PL2M's error comes
  from re-imposing a power law on the moved boxcar at every step; EA keeps the boxcar's edges as fronts, and its
  residual at few bins comes from phases where the whole boxcar sits inside one bin (held as a pivot). 1000 steps cost
  little more than 100.
\item \textbf{Growth and shrinking} (Fig.~\ref{fig:mckgrowth}b--d): a boxcar over the middle quarter of the grid with
  $\dot a$ sinusoidal (0.015\,\textmu m/Gyr, period 1\,Gyr). After one period, at 128 bins, the spread beyond the left
  edge is none for EA (and EAF), 3 bins for PL2M, 5 for PL and 15 for PC (they report 3 and 9 for PL and PC by eye;
  our threshold, $10^{-3}$ of the plateau, is stricter). The mass returns to its initial value to $3\times10^{-13}$ for
  EA at 128 and 16 bins ($7\times10^{-6}$ for PL2M, 4\,\% for PC). EA follows the analytic boxcar exactly; PL2M
  overshoots at the right edge, EAF peaks at the left one.
\end{itemize}

\begin{figure}[H]
\centering
\includegraphics[width=\textwidth]{figs/mckinnon18_growth.pdf}
\caption{After McKinnon et al.'s Figs.~10 and 13. (a) Fractional dust-mass error of the growth test at 5\,Gyr against
the number of bins, PC and PL re-implemented from their paper, PL2M, EAF and EA (default); stars: the values they quote
at 64 bins.
(b, c) The growth--shrinking cycle at 128 and 16 bins, at four times, against the analytic boxcar (black).
(d) The mass error through the cycle.}
\label{fig:mckgrowth}
\end{figure}

\subsection{The grid edge and a finite metal reservoir (their Figs.~11, 12)}
\begin{itemize}
\item \textbf{Grains growing past $a_{\max}$} (Fig.~\ref{fig:mckbound}a): grains up to 1\,\textmu m grown by
  2.4\,\textmu m with grids ending at 1, 2 and 4\,\textmu m. Rebinning the overflow into the top bin, at the top edge
  with its mass kept (McKinnon's and FIRE's choice), makes the mass grow too fast: $\times$22 by 8\,Gyr for the
  1-\textmu m grid, for PL and PL2M alike, as McKinnon et al.\ explain (radius is capped, so the number grows). Carrying
  the overflow as one population at its mean size, growing with the rest (shown with EA; the closure does not matter
  here), is within 1.8\,\% (0.5\,\% for 2\,\textmu m).
  RAMSES-CALIMA's overflow collector should therefore carry $N$ and $M$ and keep growing, not be rebinned.
\item \textbf{Decided for RAMSES (18 November): one overflow bin per composition} above $a_{\max} = 1$\,\textmu m,
  open-ended, holding $N$ and $M$ (2 scalars per composition, as one more ordinary bin):
  \begin{itemize}
  \item its grains grow, drift and collide as one population at the bin's mean size. The top bin's growth remap fills
    it; erosion returns grains to the grid;
  \item it is computed only when needed. Its charge solve, optics, drift and collisions are skipped while it holds
    less than $10^{-6}$ of the cell's dust mass (a parameter). Its share of any rate is then of order $10^{-6}$, since
    its per-gram rates are no larger than the smaller grains' to within a few. Transport always carries it. In most
    of a galaxy it stays empty and costs nothing but its two scalars;
  \item a diagnostic replaces a hard stop. Each coarse step logs its share of the dust mass and warns above 1\,\%,
    meaning $a_{\max}$ was too low for this run. An optional namelist threshold aborts the run, off by default.
    pyCALIMA has a strict mode that raises, so tests can assert it stays empty;
  \item its optics and charging tables reach about 10\,\textmu m, evaluated at its mean size; near the geometric
    limit they are cheap.
  \end{itemize}
  Why not a larger $a_{\max}$ with a stop: coagulation in dense gas can cross any edge. Sweep-up at sticking 1 gives
  $da/dt \approx \rho_{\rm d}\Delta v/(4\rho_{\rm gr})$, roughly 1\,\textmu m per Myr at $n_{\rm H}\sim10^5$\,cm$^{-3}$
  and $\Delta v\sim0.1$\,km\,s$^{-1}$, an order-of-magnitude estimate. Coreshine points to grains of about 1\,\textmu m
  in dense cores (Pagani et al.\ 2010; Steinacker et al.\ 2010). Reaching 10\,\textmu m with ordinary bins would add
  3 bins per composition, and a stop would end a large run over a few cells.
\item \textbf{A finite reservoir} (Fig.~\ref{fig:mckbound}b): dust accreting from gas-phase metals ($Z = 0.1$) with
  initial dust-to-gas ratios $10^{-2}$--$10^{-4}$, through \texttt{gsd.exchange}'s implicit gas update. With the
  physical rate, proportional to the gas-phase density, the dust mass approaches $1 + Z/D$ smoothly (McKinnon's fixed
  $\dot a$ stops abruptly); the final values match it to $2\times10^{-11}$ and $4\times10^{-7}$ (the $10^{-4}$ case is
  still converging at 4\,Gyr), with the metals conserved to round-off throughout.
\end{itemize}

\begin{figure}[H]
\centering
\includegraphics[width=\textwidth]{figs/mckinnon18_boundary.pdf}
\caption{After McKinnon et al.'s Figs.~11 and 12. (a) Dust mass of grains growing past the top of the grid (grids to
1, 2, 4\,\textmu m; marker size), with the overflow rebinned into the top bin (PL, PL2M) or carried at its mean size
(EA). (b) Growth from a finite metal reservoir with the implicit gas update; dotted lines: $1 + Z/D$.}
\label{fig:mckbound}
\end{figure}

\subsection{Shattering and coagulation (their Figs.~15, 16)}
\begin{itemize}
\item \textbf{Shattering}: their log-normal ($a_0 = 0.1$\,\textmu m, $\sigma = 0.6$) in WNM-like gas, collisions at
  3\,km/s only between grains above 0.1\,\textmu m, every collision shattering both grains into $a^{-3.3}$; 100\,Myr.
  With 8 bins (their number) the bins holding at least 1\,\% of the mass are within 15\,\% of the reference, and 2\,\%
  of the mass is misplaced in all; 7\,\% and 1\,\% with 16 bins. EA and PL2M agree to a few tenths of a per cent here,
  since collisions fill the bins smoothly. Total mass is conserved to $10^{-15}$.
\item \textbf{Coagulation}: their MRN in dense gas ($n_{\rm H} = 10^5$, $D = 0.01$), $v(a)\propto a^{1/2}$, sticking
  below 11\,m/s; 20\,Myr. The 32-bin run follows the reference except in the bin at the sticking threshold
  ($a\approx0.02$--0.03\,\textmu m), where the kernel is a step function and two nodes per bin cannot resolve it:
  24\,\% in the bin just above it (0.039--0.049\,\textmu m, where the coagulated grains pile up) and 3.5\,\% of the mass
  misplaced, from 2\,Myr on (with a reference whose 640 fine bins nest in the coarse ones; the first version's 400-bin
  reference did not, and gave 6\,\%). Before the collision solver's fix (PR \#22) this read 16\,\%; the old solver gives
  23\,\% too at a tenth of the test's error tolerance, so the 16\,\% was a step-size artefact, and the new result is
  converged in time ($10^{-4}$ and $10^{-6}$ agree). McKinnon et al.\ report ``some tension near $a\approx
  0.02$\,\textmu m'' at the same place. With 8 bins the misplaced mass is 0.3--1\,\%. EA and PL2M give the same
  numbers.
\end{itemize}

\begin{figure}[H]
\centering
\includegraphics[width=\textwidth]{figs/mckinnon18_collisions.pdf}
\caption{After McKinnon et al.'s Figs.~15 and 16. Top: shattering of a log-normal after 100\,Myr with 8 and 16 bins.
Bottom: coagulation of an MRN distribution at 2, 5, 10 and 20\,Myr with 32 and 8 bins. Lines: the fine reference;
circles: EA bin averages (thin lines: its profiles); crosses: PL2M bin averages.}
\label{fig:mckcoll}
\end{figure}

\subsection{What the tests show}
\begin{itemize}
\item EA and PL2M pass every test of McKinnon et al., and where they run the same problem as McKinnon's PL scheme
  they are more accurate, with exact number and mass, no limiter and no step limit; EA is 9 times more accurate than
  PL2M in the growth test at 8 bins and exact from 64, and exact in the growth--shrinking cycle.
\item PL2M's artefact is the shape inside partly filled bins: the power law piles grains against an edge (Figs.~9, 13).
  EA removes it by treating such bins as fronts that continue the neighbour's slope (Sect.~\ref{sec:eaf}).
\item Sharp features of a kernel (a sticking threshold) inside a bin are resolved only to that bin; a third node or
  bin edges at known thresholds would help.
\item The overflow above $a_{\max}$ must be carried and grown, not rebinned.
\end{itemize}

%======================================================================================================================
\section{What RAMSES-MCR (Diallo et al.\ 2026) offers our solver}\label{sec:diallo}

Diallo et al.\ (2026) evolve cosmic-ray protons in momentum bins with the closure PL2M uses for grains: a power law in
each bin, its slope fixed by the ratio of two moments (number and energy), tabulated against that ratio (their
Eq.~20). Their choices, and what each would do for us:
\begin{itemize}
\item \textbf{Moment-specific coefficients} (their Sect.~2.2.5, after Hopkins 2023). With one bin-centred diffusion
  coefficient, number and energy move together and the inner slope can never change; they give each moment its own
  effective coefficient. The dust analogue is the drift: with one velocity per bin, big and small grains of a bin
  cannot separate. Giving $N$ and $M$ their own velocities from the power law fixes it and keeps the bins realisable
  (Sect.~\ref{sec:pl2m}, transport). Figure~\ref{fig:drift} tests it: one 20--53\,nm bin drifting with $w\propto a$ in
  a 1-D box. One bin-centre velocity freezes $s$ and is off by 47\,\% in $N$ and $M$; the moment velocities follow the
  64-sub-bin reference to 9--12\,\% and its in-bin state to 0.04; the two Gauss nodes per cell (the default rate
  nodes) do as well (9--12\,\%, 0.035), the median nodes less so (15--23\,\%).
  The same applies to any size-dependent sink that removes grains without changing their size (molecular carbon loss,
  freeze-out onto large grains, Coulomb explosion): $N$ and $M$ need their own rates.
\item \textbf{The slope from a table} of $s$ (their $e/n$) per bin, rather than an iteration: a speed-up for the
  Fortran.
\item \textbf{Open boundaries} that give the edge bins their neighbour's slope for inflowing quantities.
\item \textbf{In-bin integrals of the loss terms}, by quadrature with a semi-implicit update, rather than bin-centre
  values: the same conclusion as our in-bin rate test (Sect.~\ref{sec:inbin}).
\end{itemize}
What it does not change: their transport in momentum is a flux through the bin edges, integrated over the power law
across the momentum swept in a step, with the step limited to a tenth of the bin's width; our characteristic remap
does the same without a step limit and is exact for a constant $\dot a$. And since the closure is the same, the
artefact in partly filled bins and the threshold resolution remain.

\begin{figure}[H]
\centering
\includegraphics[width=\textwidth]{figs/inbin_drift.pdf}
\caption{One grain bin (20--53\,nm, MRN-like start, $s = 0.14$) drifting with $w\propto a$ through a periodic box,
first-order upwind for all. Black: 64 sub-bins, each at its own velocity. Dashed: one velocity per bin (bin centre).
Blue: $N$ and $M$ with their own power-law velocities (the Diallo/Hopkins approach). Dotted: the median nodes;
dash-dot: the Gauss nodes. Right:
the in-bin state, which only the last two let evolve.}
\label{fig:drift}
\end{figure}

%======================================================================================================================
\section{How the default was chosen: variants at equal memory}\label{sec:variants}

You asked whether the defects found above can be reduced without adding memory (still two scalars per bin, $N$ and
$M$) or computational cost. A literature survey (\texttt{lit/pl2m\_like\_solvers.md}; mostly unverified leads, since it
read no full texts) fed a second, careful study that designed variants, implemented them as standalone code
(\texttt{diagnostics/gsd/pl2m\_variants.py}, \texttt{bench\_pl2m\_variants.py}, \texttt{bench\_pl2m\_figures.py}; no
existing file changed) and benchmarked them against PL2M on every test of this report plus harder stress tests. The
full report is \texttt{variants/REPORT.md}; I checked its numbers against the output files and every figure.

\emph{Dated study (17--18 November).} Its numbers are from the collision solver before PR \#22 and were not rerun: the
bench patches internals that moved into \texttt{GSDMethod} on 18--19 November, so it no longer runs. Where it overlaps
the production runs, those supersede it: the A0 numbers of EA, EAF and PL2M (Sect.~\ref{sec:binconv}; shattering at 8
bins EA 7.1\,\%, PL2M 5.7\,\%) and McKinnon's tests (Sect.~\ref{sec:mck}). The collision variants (kernel averaging,
Gauss nodes) are to be re-measured through \texttt{GSDMethod} before kernel averaging is considered as a default.

Two checks come first. With the PL2M closure the variant machinery reproduces the production remap to $3\times10^{-15}$
and the collision nodes to $3\times10^{-14}$, and a PL2M A0 rerun through it reproduces the stored run bit for bit.
Every variant keeps $N$ and $M$ exactly, stays positive and keeps $\bar m$ inside the bin.

\subsection{The variants}
\begin{table}[H]
\centering\small
\caption{Variants tested, all with $N$ and $M$ per bin. Cost: compiled closure + remap per cell and step at 8 bins,
relative to PL2M.}\label{tab:variants}
\begin{tabularx}{\textwidth}{@{}l L l@{}}
\toprule
variant & idea & cost \\
\midrule
\textbf{EA} (edge-aware) & a bin whose neighbour is (nearly) empty is a front: its profile continues the neighbour's
  power-law exponent and is cut off inside the bin (cut set by $N$, $M$); a bin with both neighbours empty is a spike at
  its mean mass; elsewhere PL2M. Emptiness threshold $\epsilon$ on the neighbour's number density per size & 1.22 \\
EAF & EA with a fixed front exponent ($-3.5$), so every shape depends on $s$ only & 1.22 \\
TC & PL2M with the exponent limited to $\pm L$; beyond, a cut-off power law & 0.8--0.9 \\
affine, exponential & other closures with the same two moments: linear in $a$ with a positivity switch; $e^{-\kappa m}$
  (maximum entropy) & 5.8, 9.3 \\
Gauss nodes & collision/rate nodes at the 2-point Gauss quadrature of the bin in radius (exact for $\int a^k dN$, $k\le3$)
  & same pairs \\
KA & kernel averaging: each node pair's rate, products and donor split averaged over $4\times4$ sub-slices & 5.5$\times$
  one collision evaluation \\
\bottomrule
\end{tabularx}
\end{table}

\subsection{Growth, edges and stress tests}
The edge-aware closure removes most of the growth error (Fig.~\ref{fig:varg}). In McKinnon's growth test at 8 bins and
100 steps the error falls from 14.2\,\% (PL2M) to 1.6\,\%, and it is exact to round-off from 64 bins, because a front
that sits inside a bin is now represented as one. In the growth--shrinking cycle the spread past the edge falls from 3
bins to none (Fig.~\ref{fig:vare}). The stress tests (Fig.~\ref{fig:vars}) show the same picture:
\begin{itemize}
\item a narrow population grown by a factor 10 to 100 over up to 5000 steps: mass error 17--20\,\% $\to$ $10^{-5}$;
\item ten growth--shrink cycles: PL2M's error grows with every cycle (L1 of the mass distribution 0.07 $\to$ 0.31), EA's
  stays at zero for a boxcar and below 0.004 for an MRN segment;
\item a cut-off at any position inside a bin: the Wasserstein distance of the mass distribution falls 3--13 times;
\item two cells mixed, 200 times: L1 error 0.12 $\to$ 0.016, surface error $-3.3\,\% \to -0.2\,\%$;
\item RAMSES-like steps ($10^2$--$10^4$\,yr, $10^3$--$10^5$ steps) on the real grids (8 silicate bins; the 11-bin carbon
  layout with $\eta$ = 2.1 and 18.8): mass error 0.8\,\% $\to$ 0.12--0.16\,\%, surface error 4\,\% $\to$ 1\,\%,
  independent of the step size.
\end{itemize}
The affine and exponential closures, with the same memory, do worse than PL2M on MRN-like profiles at 6--11 times the
cost; TC helps only marginally.

\begin{figure}[H]
\centering
\includegraphics[width=\textwidth]{variants/variants_growth.pdf}
\caption{McKinnon's growth test with the variants: (a) mass error at 5\,Gyr against the bin number (100 steps);
(b, c) against the number of steps at 8 and 32 bins. EA (red) is exact to round-off from 64 bins.}
\label{fig:varg}
\end{figure}

\begin{figure}[H]
\centering
\includegraphics[width=\textwidth]{variants/variants_edges.pdf}
\caption{How each closure represents a sharp edge. Top: a boxcar after 20 growth steps (grey: exact). Bottom: the
growth--shrinking cycle at 16 bins after half a period (dashed) and one period (solid), with McKinnon's PL for
comparison. EA keeps the edge inside the bin where it is.}
\label{fig:vare}
\end{figure}

\begin{figure}[H]
\centering
\includegraphics[width=0.85\textwidth]{variants/variants_stress.pdf}
\caption{Stress tests: each row's error metric for PL2M (open circles) and the variants. EA (red squares) is the best on
every transport test; values below $10^{-7}$ are drawn at $10^{-7}$.}
\label{fig:vars}
\end{figure}

\subsection{A0, rates and drift}
\begin{itemize}
\item \textbf{A0 at 8 bins} (worst observable over the snapshots): EA lowers accretion 8.1 $\to$ 2.8\,\%, coagulation
  5.9 $\to$ 0.3\,\%, shattering 7.8 $\to$ 4.4\,\% and sputtering 12.1 $\to$ 0.6\,\%; SN dust stays at 3.4\,\%. The worst
  over the five cases falls from 12.1 to 4.4\,\% (Fig.~\ref{fig:vara}). EA lowers accretion, coagulation and sputtering
  at every bin count (at 4 bins accretion 23 $\to$ 9\,\%, sputtering 21 $\to$ 6\,\%), but the worst case at 4 and 6
  bins is SN dust ($R_V$), which it does not change.
\item About half of that gain needs the optics and rate tables to see the EA shape: EA's states read through the
  present PL2M tables give 7.7\,\%. So EA requires front-aware tables: two-dimensional $s$ tables (in $s$ and the
  front exponent), four one-dimensional ones with EAF, or node tables, which serve any closure at about the size of the
  PL2M $s$ tables (Sect.~\ref{sec:nodetab}).
\item \textbf{The emptiness threshold must be small} ($\epsilon = 10^{-6}$): with $10^{-2}$ the very wide bins of 4--6-bin
  grids treat smooth MRN tails as fronts, and A0 errors reach 92\,\%.
\item \textbf{Rates:} with EA states and Gauss nodes in radius, the worst rate error at 8 bins is 1.0\,\% (exact in-bin
  integral) and 2.1\,\% (two nodes), against 6.2\,\% and 7.8\,\% for PL2M.
\item \textbf{Drift:} Gauss nodes bring the two-node velocities to the accuracy of the per-moment velocities (9\,\% and
  12\,\% in $N$ and $M$, from 15\,\% and 23\,\%). EA's isolated-bin spike must not be used for the drift, since it would
  freeze the in-bin segregation; there the local closure stays.
\end{itemize}

\begin{figure}[H]
\centering
\includegraphics[width=\textwidth]{variants/variants_a0_rates_drift.pdf}
\caption{A0 cases against the bin number for every variant (dashed red: EA with $\epsilon=10^{-2}$, solid red:
$10^{-6}$), the worst grain-rate error at 4 and 8 bins (filled: exact in-bin integral; open: two nodes), and the drift
test.}
\label{fig:vara}
\end{figure}

\subsection{Collisions}
No closure change helps collisions much, since collisions fill bins smoothly. Two quadrature changes do
(Fig.~\ref{fig:varc}):
\begin{itemize}
\item \textbf{Kernel averaging} fixes the sharp shattering threshold: in McKinnon's test at 8 bins the worst bin falls from
  15\,\% to 0.6\,\% and the misplaced mass from 2.1\,\% to 0.26\,\%; with the threshold swept through a bin, from
  2.1\,\% to 0.6\,\% at worst. It also improves the Golovin kernel (37 $\to$ 12\,\% at 8 bins) and cuts the sub-steps
  of the step control tenfold in threshold cases. It costs 5.5 times one collision evaluation.
\item \textbf{Gauss nodes} are mixed: much better for the Golovin kernel (37 $\to$ 4\,\% at 8 bins) and for coagulation
  at 32 bins (3.3 $\to$ 1.1\,\%), but worse for the constant kernel at 8 bins (2.3 $\to$ 12\,\%) and coagulation at 8 bins.
\item \textbf{The coagulation threshold is not fixed by anything}: sticking empties the light side of the threshold bin,
  an in-bin edge that no two-parameter description holds (16\,\% in that bin at 32 bins with the old collision solver;
  24\,\% after its fix, Sect.~\ref{sec:mck}).
\end{itemize}

\begin{figure}[H]
\centering
\includegraphics[width=\textwidth]{variants/variants_collisions.pdf}
\caption{Collision tests: McKinnon's shattering and coagulation tests (misplaced mass; small faint symbols: worst bin),
a shattering and a coagulation threshold swept through a bin, and the constant and additive kernels against their
analytic solutions.}
\label{fig:varc}
\end{figure}

\subsection{What was adopted (18--19 November)}
\begin{enumerate}
\item \textbf{EA is the default closure} of every \texttt{pycalima.gsd} operator: fronts that continue the neighbour's
  slope, $\epsilon = 10^{-6}$. EAF (a fixed front slope, $-3.5$) was the default on 18 November, to keep the tables
  one-dimensional in $s$; on flat profiles its slope puts peaks at the edges (Figs.~\ref{fig:mck9},
  \ref{fig:mckgrowth}), so EA replaced it, and its tables are kept small with node tables (Sect.~\ref{sec:nodetab}).
  The cost is about 22\,\% more in the growth operator per cell and step, no memory; EA reads the neighbouring bins of
  the same cell, which is local in RAMSES.
\item \textbf{Gauss nodes in radius are the default rate nodes}; the collisions keep the median nodes.
\item \textbf{All methods are selectable} in \texttt{pycalima.gsd.closures}. A \texttt{GSDMethod} sets the closure,
  $\epsilon$, the front exponent of EAF, the collision and rate nodes and the kernel averaging; the presets are
  \texttt{ea} (default), \texttt{eaf}, \texttt{pl2m}, \texttt{affine} and \texttt{exp}, and every operator takes a
  \texttt{method} argument. The \texttt{pl2m} preset reproduces the code before 18 November bit for bit. TC is not
  offered. Kernel averaging is an option, to be tested inside the full collision model before any default use.
\item \textbf{Checks:} 120 GSD tests pass (67 existing, 53 for the methods and node tables) and the 8 docs tests;
  through the production code the growth test and the A0 numbers match the study to four digits
  (\texttt{variants/}\allowbreak\texttt{production\_check.txt}).
\item \textbf{Accepted limits} of two degrees of freedom per bin: the coagulation-threshold bin, and the jump of the
  profile at bin edges in smooth regions.
\end{enumerate}
Every figure of this report has been redrawn with EA as the default and PL2M for comparison (and EAF where it
differs).

\subsection{Tables at a minimum for EA: per-radius node tables}\label{sec:nodetab}
With PL2M a bin's per-gram optics depend on the cell only through $s$, so RAMSES can read $\kappa_i(\lambda, s)$ from a
one-dimensional table (Sect.~\ref{sec:materials}). EA breaks this: a front bin's shape depends on its neighbour's
exponent as well as on $s$. Two-dimensional $s$ tables (in $s$ and the front exponent) would cost roughly 16 times the
PL2M tables, and even EAF's four one-dimensional tables per bin cost 4 times.

The way out uses the structure all closures share. Every reconstruction is a few power-law pieces, and the per-gram
average of a single-grain quantity $k(u)$ ($u = a/a_L$) over any piece is exact when $k$ is linear in $u$ between nodes
(\texttt{closures.bin\_weights}). So RAMSES can store $k$ at a few radii per bin, the \emph{node table}, and integrate
it over each cell's own reconstruction:
\begin{equation}
\kappa_i(\lambda) = \sum_j W_{ij}\,k_i(u_j, \lambda),
\end{equation}
where $W_{ij}$ is the exact integral of the hat function at node $u_j$ over the mass in bin $i$'s pieces. The weights
are computed once per cell and hydro step and shared by every wavelength (or RT group) and quantity. \texttt{optics.node\_table} picks, per bin, the cheapest layout (uniform or geometric in $u$, from 8 to 256
intervals) whose averages over a set of probe profiles meet a tolerance: PL2M power laws at every in-bin state and
truncated fronts with slopes from $-6$ to $+3$ cut anywhere in the bin. \texttt{optics.node\_average} evaluates the
table for any method. The same table serves EA, EAF, PL2M and the other closures.

\begin{table}[H]
\centering\small
\caption{Optics tables for the A2 grids (8 bins per material; 61 wavelengths; $C_{\rm abs}$, $C_{\rm sca}$, $C_{\rm rp}$):
stored values per wavelength for all 8 bins at a tolerance of $3\times10^{-3}$, and the worst relative error of the
per-gram $\kappa$ against a 512-node reference (\texttt{figs/node\_tables.txt}). The $s$ tables serve PL2M only; EAF's
four component tables are 4 times the $s$ tables; EA's two-dimensional $s$ tables are an estimate.}\label{tab:nodetab}
\begin{tabularx}{\textwidth}{@{}l r r L@{}}
\toprule
tables & carbon & silicate & closures served, worst error \\
\midrule
PL2M $s$ tables (A2) & 358 & 298 & PL2M; 2.8--3.2$\times10^{-3}$ \\
EAF component $s$ tables & 1432 & 1192 & EAF \\
EA as 2-D $s$ tables (estimate) & $\sim$5700 & $\sim$4800 & EA \\
\textbf{node tables} (tol $3\times10^{-3}$) & \textbf{336} & \textbf{304} & \textbf{any}; $2.9\times10^{-3}$ (carbon);
  $2.1\times10^{-3}$ on EA fronts, $4.0\times10^{-3}$ on random states (silicate) \\
node tables, best layout per bin & 298 & 318 & any; $\le3\times10^{-3}$ \\
\bottomrule
\end{tabularx}
\end{table}

\begin{figure}[H]
\centering
\includegraphics[width=\textwidth]{figs/node_tables.pdf}
\caption{Node tables against $s$ tables: worst relative error of the per-gram $\kappa$ against the stored values per
wavelength (all 8 bins), for uniform, geometric and adaptive node layouts and the production builder, on PL2M profiles
(open) and EA states with fronts (filled); stars: the A2 $s$ tables, which serve PL2M only.}
\label{fig:nodetab}
\end{figure}

\begin{itemize}
\item \textbf{Memory:} node tables cost about as much as PL2M's $s$ tables (carbon 336 against 358 values per
  wavelength; silicate 304 against 298) and serve EA exactly, against 4 times for EAF's tables and an estimated 16
  times for two-dimensional EA tables. The band-gap axis of the M1 bins multiplies node and $s$ tables alike.
\item \textbf{Cost:} the weights take 6--11\,\textmu s per cell and hydro step for 8 bins (compiled, one thread),
  shared by every group and quantity; each group then costs one dot product per bin. With PL2M's $s$ tables this was a
  lookup, so the optics side of EA adds about the cost of one closure evaluation per hydro step.
\item \textbf{Consistency:} the same per-radius approach serves the charge kernels at the Gauss nodes
  (Sect.~\ref{sec:inbin}) and the drift velocities, so optics, rates and transport all integrate the cell's own
  reconstruction.
\item \textbf{Where they are used:} in pyCALIMA, for exact evaluation, and in the post-processing export (A6). For the RT
  cross-sections in RAMSES's cool step even these weights are too dear; RAMSES uses five pseudo-bins per bin instead
  (Sect.~\ref{sec:rtgroups}).
\end{itemize}

\subsection{RT group cross-sections in the cool step: five pseudo-bins per bin}\label{sec:rtgroups}

\paragraph{How RAMSES does it today.} RAMSES-RT tabulates each dust bin's per-gram group cross-sections against stellar
age and metallicity at start-up, one copy per node. These are absorption (csa), scattering (css) and radiation pressure
(csr), stored in \texttt{init\_SED\_table}. Every few coarse steps it weights them by the stars' luminosities
(\texttt{update\_SED\_group\_props}), and the cool step copies the result into each cell. The IR group (0.1--1\,eV) is
the exception: on every substep its cross-sections are replaced by Planck and Rosseland means at the cell's radiation
temperature. While each bin had one fixed size distribution, its cross-section was the same in every cell. With two
moments it depends on the cell's in-bin profile.

\paragraph{Why not EA exactly.} The exact EA average costs 4.0\,\textmu s per cell-call for 8 carbon and 8 silicate bins
(compiled, one thread). It takes $s$ tables where a bin is pure base, where EA equals PL2M, and node integrals over the
fronts. That is a large addition to a cool step of 30--80\,\textmu s, for numbers the RT uses only as absorption and
radiation pressure.

\paragraph{The scheme.} The cool step ignores EA's fronts and uses the PL2M profile at the bin's $s$. Each GSD bin is
exported as five \emph{pseudo-bins}: the PL2M per-gram $\sigma(\lambda)$, with its Planck and Rosseland tables, at the
fixed states
\begin{equation}
s_k = \tfrac12\left[1 - \cos(\pi k/4)\right] = 0,\ 0.15,\ 0.5,\ 0.85,\ 1 .
\end{equation}
These are Chebyshev points, which cluster near the bin edges. RAMSES treats every pseudo-bin as an ordinary dust bin:
in the SED table, in the star polling and in the IR means, all unchanged. This works because the group average, the
luminosity weighting and the blend below are all linear in $\sigma$. In each cell-call every bin finds its bracket and
weight once,
\begin{equation}
\kappa_{i,g} = (1-t_i)\,\kappa^{(k)}_{i,g} + t_i\,\kappa^{(k+1)}_{i,g}, \qquad t_i = \frac{s_i - s_k}{s_{k+1}-s_k},
\end{equation}
and keeps them over the substeps, while the masses still update. The IR substep uses one $T_{\rm rad}$ index for all
bins and blends the two bracketing pseudo-bins. The shielding bands, and csrat and la when they are on, take the same
blend. RAMSES copies css into the cell but none of its rates read it, so css is dropped.

\begin{table}[H]
\centering\small
\caption{Per-gram RT group cross-sections against the exact EA in-bin average: $|\mathrm{scheme}/\mathrm{EA}-1|$ in
\%, worst over the states, for the composition total $\sum_i M_i\kappa_i$ (what the RT absorption sees); in brackets,
per bin (bins with $\ge$\,1\,\% of the mass). Groups 2--8: RTZ's groups, photon-weighted over a $5\times10^4$\,K
blackbody, absorption and radiation pressure. IR: Planck absorption and Rosseland radiation pressure at
$T_{\rm rad} = 10$--300\,K. A0: the EA 8-bin end states of the five cases. Stress: fronts cut anywhere in a bin, narrow
segments and random states (135). $\times$: the ratio, where the error exceeds 100\,\%.
\texttt{diagnostics/gsd/check\_rt\_group\_optics.py}, \texttt{figs/rt\_group\_optics.txt}; A0 columns rerun on 22
November on the new A0 states.}\label{tab:rtgroups}
\begin{tabular}{@{}l rr rr rr rr@{}}
\toprule
 & \multicolumn{4}{c}{A0 end states} & \multicolumn{4}{c}{stress states} \\
\cmidrule(lr){2-5}\cmidrule(l){6-9}
 & \multicolumn{2}{c}{groups 2--8} & \multicolumn{2}{c}{IR} & \multicolumn{2}{c}{groups 2--8} & \multicolumn{2}{c}{IR} \\
scheme & carbon & silicate & carbon & silicate & carbon & silicate & carbon & silicate \\
\midrule
one value per bin ($s=\tfrac12$) & 25 (57) & 27 (66) & 10 & 2.8 & 44 & $\times2.4$ & $\times66$ & 19 \\
grain at the mean mass & 7.9 (11) & 9.5 (25) & 3.5 & 0.3 & 18 & 9.8 & 46 & 5.0 \\
3-point polynomial in $s$ & 28 (40) & 31 (40) & 7.3 & 0.2 & 38 & 33 & $\times5.8$ & 6.0 \\
3 pseudo-bins & 36 (50) & 42 (53) & 9.2 & 0.3 & 53 & 47 & $\times5.0$ & 7.9 \\
\textbf{5 pseudo-bins (adopted)} & \textbf{4.1 (15)} & \textbf{6.1 (14)} & \textbf{1.4} & \textbf{0.1} & \textbf{17} &
  \textbf{12} & $\boldsymbol{\times7.8}$ & \textbf{3.9} \\
9 pseudo-bins & 1.8 (5.9) & 2.1 (5.6) & 0.9 & 0.2 & 15 & 12 & $\times5.0$ & 4.1 \\
PL2M at $s$, exact & 1.8 (3.1) & 3.6 (6.1) & 0.8 & 0.2 & 14 & 18 & $\times5.2$ & 4.4 \\
\bottomrule
\end{tabular}
\end{table}

\begin{itemize}
\item \textbf{Accuracy} (Table~\ref{tab:rtgroups}): on the A0 states five pseudo-bins are within 4--6\,\% of EA on the
  group totals and within 1.4\,\% on the IR means (1.3\,\% before the 22 November rerun). That is the error you accepted on 18 November. Nine pseudo-bins
  would reach PL2M's own floor (about 2\,\%).
\item \textbf{What it gives up:} the fronts. Where a population ends inside a bin, the RT optics see the PL2M profile,
  not EA's: up to 12--17\,\% on the group totals in the stress states. On narrow carbon segments the IR opacity is
  5--8 times too large, against 4--5 times for PL2M itself. The dust solver, the rates and the drift still use EA.
\item \textbf{Cost} (\texttt{bench\_rt\_group\_optics.py}): 0.19\,\textmu s per cell-call against 0.11 for today's
  copy; 9 pseudo-bins take 0.21, PL2M's full $s$ tables 0.30, EA exactly 4.0. The IR substep costs 0.03\,\textmu s, as
  today.
\item \textbf{Memory:} the SED table's dust columns (247 ages, 8 metallicities, 8 groups, real*8, one copy per node)
  grow from 6.1 to 20.2\,MB for 16 bins, or 36\,MB with nine pseudo-bins. The per-bin $\sigma(\lambda)$ tables stay
  small.
\item \textbf{The polling stays.} Across BPASS constant-SFR mixes, a bin's per-gram group values vary by 30--53\,\% in
  group 2 (1--5.6\,eV), by 5--13\,\% in groups 7--8 and by at most 10\,\% elsewhere. Across single bursts the variation
  reaches 95\,\%. A fixed spectrum would cost more than the in-bin error.
\item \textbf{Open for the hydrogen-tracking bins (M1):} the band-gap axis multiplies the pseudo-bins, giving a bilinear
  blend in $s$ and $\Eg$. Five $s$ points times about six $\Eg$ nodes on the 8 bins below 53\,nm would take about
  60\,MB of the SED table per node. Fewer $\Eg$ nodes or a coarser age axis have to be measured before B6/B7.
\end{itemize}

%======================================================================================================================
\section{How many bins, and what the bins must reproduce}\label{sec:bins}

\subsection{Optical observables (A0)}\label{sec:binconv}
\texttt{diagnostics/gsd/check\_bin\_convergence.py} runs five fast test cases against reference solutions. Two of them,
accretion and sputtering, have exact solutions. The metrics are observables, not the distribution itself: mass, slope,
mean and width in $\ln a$, $A_V$, $R_V$, the bump, the UV slope and the extinction curve.
\begin{itemize}
\item The two-moment schemes need about half the bins of a one-moment scheme for the same accuracy
  (Fig.~\ref{fig:binconv}).
\item With 8 bins per composition, the worst observable errors with EA (PL2M in brackets) are:
  \begin{itemize}
  \item SN injection into diffuse gas: 3.1\,\% (3.1\,\%);
  \item accretion: 2.8\,\% (8.1\,\%);
  \item coagulation: 0.3\,\% (5.8\,\%);
  \item shattering: 7.1\,\% (5.7\,\%), on the graphite mean size;
  \item sputtering: 0.6\,\% (12.1\,\%, on the silicate width).
  \end{itemize}
  The worst over the five cases is 7.1\,\% with EA (7.2\,\% with EAF) against 12.1\,\% with PL2M.
\item With 6 bins the worst is 19\,\% for EA and 18\,\% for PL2M (shattering); with 12 bins 3.9\,\% and 4.9\,\%; with 32
  bins 0.6\,\% and 2.2\,\%.
\item \textbf{Rerun on 22 November} with the fixed collision solver (PR \#22). Before it, shattering read 4.4\,\% (EA)
  and 7.8\,\% (PL2M) at 8 bins, SN dust 3.4\,\%, and the worst over the cases 4.4\,\% with EA; at 6, 12 and 32 bins
  17, 4.6 and 1.7\,\% (EA). The reference moved by only 0.2\,\%; the coarse runs moved (EA's 8-bin graphite mean
  size by 6.8\,\%) and are now converged in time (errmax $10^{-5}$, twelve times the substeps, agrees to 0.01\,\%). So
  the old 4.4\,\% was partly a cancellation with the old step bias; shattering now converges steadily from 8 to 32 bins
  (Fig.~\ref{fig:binconv}). At 8 bins EA misses 5\,\% only in shattering; 12 bins would meet it in every case. You
  kept 8 bins (22 November), to hold the operations and memory down.
\end{itemize}
You chose 8 bins after comparing memory per scalar: each scalar costs 16\,B per cell, and each TVA species another
24\,B. EA does not change the memory.

\begin{figure}[H]
\centering
\includegraphics[width=\textwidth]{figs/gsd_bin_convergence.png}
\caption{A0: worst relative error over the observables against the number of bins per composition (top) and against
the scalars per cell (bottom), for EA (the default, blue), EAF (purple), PL2M (orange) and two one-moment schemes. The
dashed line marks 5\,\%. Results in \texttt{results/binconv\_v4/} (PL2M, one-moment) and
\texttt{gsd\_overview/variants/a0\_prod/} (EA, EAF).}
\label{fig:binconv}
\end{figure}

\subsection{The in-bin state beyond optics}\label{sec:inbin}
Today's RAMSES-CALIMA builds its charging, photoelectric-heating, recombination-cooling, sputtering, gas--grain and
H$_2$ tables at one radius per bin, the bin centre. With 8 bins the grain bins are wide (a factor 2.66 in radius), and a
bin's rate per gram depends on where its mass sits, exactly as its opacity does.

The rates RAMSES needs from the grains are sums over grains of single-grain rates that depend only on the grain's size
and the local conditions; a grain's charge distribution depends on the gas and the radiation, not on the other grains.
So $\Gamma_i = \int_{I_i}\Gamma(a)\,dN$, and per gram $\Gamma_i/M_i = \langle\Gamma/m\rangle_i(s)$, Eq.~(\ref{eq:kavg}).
The same holds for recombination cooling, gas--grain energy exchange, H$_2$ formation, accretion and sputtering.
\texttt{diagnostics/gsd/check\_inbin\_rates.py} compares these ways to obtain a bin's rate at state $s$, for the
carbon (a-C, \texttt{ac} charging) and silicate grain bins in the WNM, CNM, a PDR and dense gas:
\begin{itemize}
\item \textbf{exact}: the integral over the bin's power law, Eq.~(\ref{eq:kavg}); what an $s$ axis on the tables gives;
\item \textbf{bin centre}: every grain at the centre (today);
\item \textbf{mean mass}: every grain at $\bar m_i$, one charge solve per bin. It is exact only if the rate per gram is
  the same at every size;
\item \textbf{area node}: one charge solve per bin at
  \begin{equation}
  a_* = \frac{\int a^3\,dN}{\int a^2\,dN}, \qquad \Gamma_i = \frac{\int a^2\,dN}{a_*^2}\,\Gamma(a_*),
  \label{eq:areanode}
  \end{equation}
  the one-point Gauss rule of the area measure $a^2\,dN$. The bin's total grain area is exact, and so is the rate
  whenever $\Gamma(a)/a^2$ is linear in $a$ across the bin. The charge-dependent rates scale roughly as the grain area,
  which the mean mass gets wrong. Both moments come from the closure at no cost, and per gram the rule is simply the
  single-grain value at $a_*$;
\item \textbf{median nodes}: the two collision nodes of Sect.~\ref{sec:coll};
\item \textbf{Gauss nodes}: the 2-point Gauss rule of the bin in radius, exact for $\int a^k dN$ with $k\le3$ (the
  default rate nodes since 18 November).
\end{itemize}
The profile here is the power law of a single bin: the question is how to evaluate a rate over a given in-bin shape.

\begin{figure}[H]
\centering
\includegraphics[width=\textwidth]{figs/inbin_rates.pdf}
\caption{Grain rates of a single bin against its in-bin state. (a) Photoelectric heating per gram of three a-C bins in
the CNM, relative to the exact profile average: bin centre (solid), mean mass (dotted), the area node (thick
dash-dot-dot, adopted), median nodes (dashed), Gauss nodes (dash-dot). (b--d) The
worst error over $s\in[0,1]$ per bin, for photoelectric heating, recombination cooling and grain surface (H$_2$
formation, accretion, sputtering), carbon (filled) and silicate (open), in four cells.}
\label{fig:inbin}
\end{figure}

Results (Fig.~\ref{fig:inbin}; \texttt{figs/inbin\_rates.txt}):
\begin{itemize}
\item Inside one bin the rate per gram changes by a factor 1.1--8 between $s=0$ and $s=1$.
\item The bin centre is off by up to a factor 2.8 at the extremes of $s$, although within 9\,\% at the MRN slope.
\item The mean mass is within 10\,\% in 95\,\% of the cases (worst 18\,\%).
\item The area node, with the same single solve, is within 1.4\,\% in half the cases and 7.0\,\% in 95\,\% of them;
  its worst, 9.7\,\%, is the largest carbon bins (0.14--1\,\textmu m) in cold gas. It is exact for the surface.
\item The median nodes are within 5\,\% (worst 9.1\,\%, the largest carbon bin in the WNM).
\item The Gauss nodes are within 1.2\,\% (worst 1.6\,\%), at the median nodes' cost of two solves per bin: they
  became pyCALIMA's default rate nodes, and stay so for size-dependent sinks, which must move $N$ and $M$ apart.
\item For surface processes the exact value is analytic: per gram $\propto\langle a^2\rangle/\langle a^3\rangle$,
  known from $\beta$ at no cost.
\item \textbf{Rerun on 22 November} with the charge solve ported from RAMSES (PR \#19; \texttt{gsd.charging} now on LG5).
  Before: area node 1.6\,\% / 7.8\,\% (median / 95th percentile), mean mass worst 19\,\%, median nodes 6\,\% (worst
  10\,\%), Gauss nodes 1.8\,\% (worst 4.2\,\%). The changes sit in the large carbon bins (53\,nm--1\,\textmu m) of the
  PDR cell, whose charge distributions are wide; e.g.\ for 53--141\,nm the area node's heating error went from 7.9 to
  1.1\,\% and the Gauss nodes' from 4.2 to 0.4\,\%. The worst area-node case (9.7\,\%) is unchanged.
\end{itemize}

\paragraph{Decided for RAMSES (18 November).} By kind of rate:
\begin{itemize}
\item \textbf{Charge-dependent rates} (photoelectric heating, recombination cooling, Coulomb factors, grain-assisted
  recombination): one charge solve per bin at its area node, Eq.~(\ref{eq:areanode}), from the EA reconstruction. That
  keeps today's number of solves; two Gauss nodes would double the charge solve, the dearest part of the dust step,
  for at most 3 points less error at 8 bins (Sect.~\ref{sec:rateconv}). $a_*$ moves inside the bin with $s$, so the
  charge kernels are still needed at a few radii per bin, interpolated in $\ln a$. That multiplies the kernel memory
  by the radii per bin, which the planned coarsening (real*4, a cut potential range) has to absorb.
\item \textbf{Size-dependent sinks} (molecular carbon loss, Coulomb explosion, freeze-out): two nodes, per moment
  (Sect.~\ref{sec:diallo}), since one node cannot move $N$ and $M$ apart.
\item \textbf{Surface processes} (H$_2$ formation, neutral accretion, gas--grain collisional cooling, sputtering mass
  loss): the exact second moment from $\beta$; no table changes.
\item \textbf{Optics:} in RAMSES's cool step, five PL2M pseudo-bins per bin, blended in $s$
  (Sect.~\ref{sec:rtgroups}); elsewhere node tables (Sect.~\ref{sec:nodetab}), the same per-radius approach as the
  charge kernels.
\item \textbf{Molecular bins:} their tables average PAHdb molecules with fixed in-bin weights. The same treatment
  (two half-bin tables weighted by the two nodes) is possible; the bins are narrow (a factor 1.28 in radius), so the
  effect should be smaller, but it has to be measured.
\item The alternative, an $s$ axis on every rate table, is exact but costs 30--70 times the memory of each table; it
  suits only small tables.
\end{itemize}

\subsection{Convergence of the thermal and chemical rates}\label{sec:rateconv}
A0 tested the observables of the distribution (mass, shape, extinction). RAMSES also needs the rates the grains impose
on the gas. A first look, \texttt{diagnostics/gsd/check\_rates\_convergence.py}, takes the final distributions of the
A0 runs (graphite and silicate, 2--32 bins and the reference), for EA and PL2M, and computes the population's
photoelectric heating, recombination cooling and total grain surface in the four cells above, each run with its own
closure.

\begin{figure}[H]
\centering
\includegraphics[width=\textwidth]{figs/rates_convergence.pdf}
\caption{Population rates against the number of bins, worst over the five A0 cases and the four cells: EA (blue,
default) with the exact in-bin integral and with two Gauss nodes per bin; EA with one charge solve per bin, at the
area node (thick green, adopted for RAMSES) or at the mean mass (purple); PL2M (orange) with the exact integral and
the median nodes; and bin-centre tables (red, today's RAMSES). Dotted lines: 8 bins and 5\,\%.}
\label{fig:rateconv}
\end{figure}

Results at 8 bins (Fig.~\ref{fig:rateconv}; \texttt{figs/rates\_convergence.txt}):
\begin{itemize}
\item EA runs, exact integral over the EA profile: within 4.4\,\% for heating, 4.6\,\% for cooling and 6.2\,\% for the
  surface; with the Gauss rate nodes 5.4--7.1\,\%;
\item EA runs, one solve per bin: at the area node 8.8\,\% for heating, 7.3\,\% for cooling, 5.8\,\% for the surface;
  at the mean mass 13--14\,\%. The worst case is shattering for both. In the other four cases the area node is within
  0.7--2.0\,\%, against 0.3--1.0\,\% for the Gauss nodes and 4.5--9\,\% for the mean mass (Table~\ref{tab:areanode});
\item PL2M runs, exact integral: within 6.1\,\%; with the median nodes 7.0--8.1\,\%;
\item bin-centre tables: off by 42--50\,\%; 5\,\% would need 32 bins.
\item \textbf{Rerun on 22 November} (PRs \#19 and \#22): every rise comes from the shattering case, whose 8-bin state
  changed with the collision fix (Sect.~\ref{sec:binconv}). Before: EA exact 3.6 / 3.9 / 5.4\,\%, Gauss nodes
  4.7--6.4\,\%, area node 7.9 / 6.6 / 5.0\,\%, mean mass 12--13\,\%, PL2M exact 6.0\,\%, median nodes 6.7--7.8\,\%, bin
  centre 41--49\,\%. The charge solve alone moves the other four cases by at most 0.3 points.
\end{itemize}
\begin{table}[H]
\centering\small
\caption{Heating / cooling of the whole population at 8 bins, EA runs: $|$run/ref$-1|$ in \%, worst over the four
cells, against the high-resolution reference integrated exactly (\texttt{figs/rates\_convergence.txt}). Correcting
the area node with the rate's size slope from the neighbouring bins' solves gains less than one point (shattering
8.1 / 6.8\,\%), so it is not used. Rerun on 22 November; the shattering column read 3.6 / 3.9, 4.7 / 6.4, 7.9 / 6.6 and
11.8 / 13.3 before the collision fix.}\label{tab:areanode}
\begin{tabular}{@{}l c rrrrr@{}}
\toprule
way & solves per bin & SN diffuse & accretion & coagulation & sputtering & shattering \\
\midrule
exact integral & many & 1.1 / 0.7 & 0.8 / 0.7 & 0.3 / 0.1 & 0.2 / 0.2 & 4.4 / 4.6 \\
2 Gauss nodes & 2 & 1.0 / 0.8 & 1.0 / 0.9 & 0.7 / 0.5 & 0.3 / 0.3 & 5.4 / 7.1 \\
\textbf{area node (adopted)} & \textbf{1} & \textbf{2.0 / 0.9} & \textbf{1.2 / 0.9} & \textbf{1.6 / 0.7} &
  \textbf{1.0 / 1.1} & \textbf{8.8 / 7.3} \\
mean-mass grain & 1 & 7.2 / 7.5 & 4.5 / 4.6 & 8.9 / 8.1 & 8.8 / 8.5 & 12.7 / 14.1 \\
\bottomrule
\end{tabular}
\end{table}

So 8 bins carry the thermal rates as well as the optics, \emph{provided} the rates are taken over the in-bin profile;
EA gains about a quarter over PL2M for heating and an eighth for cooling here (4.4 / 4.6\,\% against 6.1 / 5.3\,\%). (Gauss nodes in the collisions too, an option, would
bring the EA rate errors to 1--2\,\%, Sect.~\ref{sec:variants}.)
Caveats: graphite (the A0 material), final states only, fixed cells rather than each case's own gas, no molecular bins,
and no feedback of the rates on the gas.

\paragraph{The full test} (part of A5, once the chemistry of every species is final). Rerun the A0 cases and one-zone
PDR and diffuse-ISM runs with 6, 8, 12 and 16 bins and a reference, with the rates coupled to the gas, and score:
\begin{itemize}
\item photoelectric heating, recombination cooling, gas--grain collisional energy exchange;
\item H$_2$ formation on grains and on the molecular bins;
\item grain-assisted recombination of H$^+$, C$^+$ and other ions, and the electron density;
\item accretion (depletion) rates per element; carbon loss of the molecular bins;
\item dust temperature and IR emission;
\item the resulting gas state in thermal and chemical equilibrium ($T$, $x_e$, $f_{\rm H_2}$).
\end{itemize}
Pass: 5\,\% on each rate and on the equilibrium gas state, at 10, 30 and 100\,\% of every run, with the chosen in-bin
treatment, and on path-integrated rates (the path budgets used for the PAH tables).

%======================================================================================================================
\section{Carbon materials, grid and optics (A1, A2, A3a, step 1)}\label{sec:materials}

\paragraph{Why amorphous carbon, not graphite.} The arguments, each with what its sources actually support
(verified on 17 November; \texttt{lit/graphite\_and\_pah\_size.md}):
\begin{itemize}
\item \textbf{Stellar carbon dust is amorphous and aromatic-rich.} Gas-phase condensation experiments make soot built
  from PAHs below about 1700\,K (J\"ager et al.\ 2009). Presolar graphite exists but makes up about $10^{-6}$ of
  meteorite mass and is biased towards survivors (Draine 2003, Table~3). Jones et al.\ (2017) find ``no evidence to
  support it [graphite] as the most abundant interstellar carbon dust material''. THEMIS adopts 1.6\,g\,cm$^{-3}$ for
  the a-C nano-particles and 1.51--1.57 for the large grains (Jones et al.\ 2017, Table~1); we use 1.6.
\item \textbf{The 2175\,\AA\ bump.} Its width varies between sightlines at nearly constant central wavelength, and
  ``the graphite hypothesis does not appear to have a natural way to accommodate'' that (Draine 2003, Sect.~3.3,
  citing Draine \& Malhotra 1993).
\item \textbf{The 3.4\,\textmu m band.} Diffuse sightlines show aliphatic C--H absorption (Sandford et al.\ 1991, via
  Draine 2003, Sect.~3.4), which pure graphite cannot produce. This rules out pure graphite, not a graphite plus
  hydrocarbon mixture (which Draine 2003 himself suggested).
\item \textbf{Far-IR emission per extinction.} The graphite--silicate model of Draine \& Li (2007), fitted to Planck,
  predicts $A_V$ about twice that measured towards quasars (Planck Collaboration Int.\ XXIX 2016); dust models with
  laboratory a-C and silicate optics reproduce the emission per extinction best (Fanciullo et al.\ 2015).
\item \textbf{Polarisation} constrains the grains' structure, not graphite as such: the unpolarised 3.4\,\textmu m band
  where silicates are polarised rules out hydrocarbon mantles on aligned silicates (Chiar et al.\ 2006; Jones et al.\
  2017, Sect.~5.1), and the submm-to-optical polarisation ratio led Hensley \& Draine (2023) to one composite material
  for the large grains.
\item \textbf{Counterpoints.} Graphite reproduces the bump strength with an acceptable carbon abundance (Draine 2003);
  Draine \& Li (2007) and Hensley \& Draine (2023) keep graphitic optics in their carbonaceous grains; the bump carrier
  is still unidentified. \emph{Correction to the first version of this report:} Hensley \& Draine (2023) did not drop
  graphite.
\end{itemize}
Graphite therefore stays as the control material. THEMIS and DustEM serve as guidance and benchmarks, not as fixed
choices: where newer physics exists it is the default, and the DustEM version is kept for comparison.

\paragraph{Why molecular properties stop at $\NC\approx400$.} This is a modelling choice with support, not an
observational requirement, and it is a run-time knob.
\begin{itemize}
\item \textbf{The bands come from small molecules.} Size diagnostics give a mean of $\NC\approx50$ at the NGC\,7023 PDR
  surface and 70 in its cavity (Croiset et al.\ 2016), and PAHs below about 50--60 C do not survive in PDRs or the
  diffuse ISM (Montillaud et al.\ 2013; Murga et al.\ 2022). A single UV photon heats a small molecule far more than a
  large one, so the 3.3\,\textmu m band is radiated by $\NC\lesssim100$ and the 7.7\,\textmu m band by $\NC\sim10^2$
  (Draine et al.\ 2021).
\item \textbf{Large compact PAHs stop matching the bands.} Above about 150 C the computed 6.2\,\textmu m band broadens
  and shifts and the 8.6\,\textmu m band overtakes the 7.7, so such PAHs ``are unlikely to be the dominant species''
  (Ricca et al.\ 2012). Their spectra do not converge to a size-independent molecular limit; the molecular detail
  there is simply not constrained by the observed bands.
\item \textbf{What is seen at a few hundred C is not a molecule.} Evaporating very small grains have been interpreted
  as PAH clusters of at least 400 C (Rapacioli et al.\ 2005, 2006; Bern\'e et al.\ 2007) of mixed aromatic and
  aliphatic composition (Pilleri et al.\ 2015). In our model such carriers are a-C(:H) nano-grains with Jones' (2012)
  size-dependent optical constants, which reach down to a few tens of C.
\item \textbf{Nothing is switched off above 400 C.} Those grains still absorb, heat the gas, charge and emit, through
  the a-C optics, the \texttt{ac} charging and THEMIS-like emission; only the molecular detail (PAHdb spectra and the
  multi-photon chemistry) stops.
\item \textbf{Counterpoints.} In Draine et al.'s (2021) model the 17\,\textmu m band is radiated by $\NC\gtrsim10^3$;
  Draine \& Li (2007) keep PAH cross sections up to $a = 50$\,\AA\ ($\NC\approx6\times10^4$, a transition they call
  ``entirely ad hoc''), and Hensley \& Draine (2023) up to about 1000 C.
\end{itemize}
So $\NC^{\max}$ stays a run-time choice (through the table headers); the seam test checks that the two sides agree at
the switch, and the extended band test (open) sets the default: the 3.3, 7.7, 11.2 and 17\,\textmu m band ratios with
$\NC^{\max} = 400$ against a model that keeps molecules to about 1000 C. With 3--4 two-moment bins over $\NC$
20--400 the 3.3\,\textmu m-based ratios are within 2--7\,\%.

The grid: a material registry in \texttt{gsd/materials.py}, and a graded grid in \texttt{gsd/carbon.py}, 4 molecular
bins over $\NC$ 20--400 ($\eta \approx 2.1$, 0.39--1.06\,nm), then 7 grain bins ($\eta \approx 19$) up to 1\,\textmu m.

\paragraph{Optics.}
\begin{itemize}
\item THEMIS CM20 is read from DustEM's tables, with Henke data below 40\,nm, and reproduces the J13 extinction to 0.9\,\%.
\item Large grains are recomputed with a coated Mie code (Jones' 100\,nm $n,k$), which matches DustEM's $Q_{\rm CM20}$
  to 0.1\,\% in the continuum. DustEM's own table is broken for $a \gtrsim 0.9$\,\textmu m below 0.1\,\textmu m.
\item The a-C(:H) $n,k$ are Jones (2012), with sizes from 0.33 to 100\,nm and $\Eg$ from $-0.1$ to 2.67\,eV.
\item Between Jones' tabulated radii the far-IR goes wrong unless the size-limited band gap is followed. The fit
  $E_{\rm lim} = 0.75\,(1\,{\rm nm}/a)^{1.06}$\,eV reproduces $Q_{\rm CM20}$ between 1 and 10\,nm (Fig.~\ref{fig:gap}).
\end{itemize}

\paragraph{Per-bin tables.} A bin's opacity depends on where its mass sits inside it, i.e.\ on its in-bin state $s$
(Sect.~\ref{sec:s}), so a fixed in-bin profile does not work.
\begin{itemize}
\item We tabulate $\kappa_i(\lambda, s)$ per gram on an adaptive $s$ grid: 31--77 points per carbon bin, 32--47 per
  silicate bin (Fig.~\ref{fig:optics}).
\item Errors: $\kappa$ 0.32\,\%, observables 0.045\,\%, emitted power 0.46\,\%, dust temperature 0.07\,\%.
\item A uniform $s$ grid fails: the sub-nanometre THEMIS grains make bin 1's near-IR opacity vary 300-fold with $s$.
\item For tracked hydrogen (M1) the covalent bins get an extra axis, the bulk band gap $\Eg$, with 39--48 points per bin.
\item With the EA closure (the default since 19 November) the same per-gram data are stored as node tables: values at a
  few radii per bin, integrated over each cell's reconstruction (Sect.~\ref{sec:nodetab}).
\item The tables keep every wavelength, so RAMSES can average them over any group set. That costs 1.4\,MB as 8-group
  tables, or 42\,MB per wavelength at start-up.
\end{itemize}

\begin{figure}[H]
\centering
\includegraphics[width=0.74\textwidth]{figs/gsd_optics_tables.png}
\caption{A2: absorption opacity per gram for the smallest and a middle bin, carbon and silicate, at five in-bin states
$s$. In bin 1, carbon's near-IR opacity changes by orders of magnitude with $s$; this is why the tables carry $s$.}
\label{fig:optics}
\end{figure}

\begin{figure}[H]
\centering
\includegraphics[width=0.74\textwidth]{figs/size_limited_gap.png}
\caption{Step 1: the size-limited band gap. (a) The fitted gap (green) against DustEM's charging law. (b, c) Far-IR
absorption of a-C between Jones' tabulated radii: plain interpolation (blue) misses DustEM by up to 60\,\%, the fitted
gap stays within 15\,\%. (d) The electron affinity that follows. You approved this choice for the \texttt{ac}
charging too.}
\label{fig:gap}
\end{figure}

%======================================================================================================================
\section{Charging, heat capacity and the PAH/grain seam (A3b)}

\paragraph{Three charging materials} for carbon grains, in \texttt{wdb06\_yields} and \texttt{rt\_group\_charging}:
\begin{itemize}
\item \texttt{ac}, the default. The ionisation potential follows WD01 with work function $W = 4.75$\,eV and radii at
  $\rho = 1.6$. The photoelectric yield is the extended Kimura model, eK16 (Hrodmarsson+25), with the size-limited gap
  from 1\,nm.
\item \texttt{ach}, for a-C:H: bulk ionisation energy 4.6\,eV (Schafer+93), with Jones' $\Eg = 2.5$\,eV optics.
\item \texttt{cm20}, a port of DustEM's model, kept as the benchmark. It matches DustEM's $\langle Z \rangle$ to
  0.1\,$\sigma$ and the photoelectric current to $<$\,1\,\% (with WDB06 photo-secondaries off). Along the way:
  \begin{itemize}
  \item DustEM's shipped \texttt{ZDIST.RES} turned out wrong above about 0.5\,\textmu m; a fresh build agrees with us;
  \item a threshold in the gas-secondary yield was fixed (energy conservation).
  \end{itemize}
\end{itemize}
Intermediate H/C uses \texttt{acx} materials at 5 band-gap nodes per bin, which keep heating and cooling within 1\,\%.
Interpolating between the two end members alone is 7--9\,\% off. Against DustEM, \texttt{ac} heats the smallest grains
more efficiently (Fig.~\ref{fig:charging}): the heating per absorbed power is $\times$1.4--3.4 at 0.4--3\,nm, from the
lower thresholds and the eK16 yield, and agrees within 6\,\% from 9\,nm to 1\,\textmu m.

\paragraph{Heat capacity.} A harmonic per-atom model (\texttt{carbon\_heat\_capacity.py}): sp$^2$ and sp$^3$ modes plus
CH$_2$ modes per H, taking $\Eg$. DustEM's $C_{\rm CM20}$ is 1.5--1.7 times larger at 5000\,K (a liquid-polyethylene
extrapolation) and is kept as the benchmark.

\begin{figure}[H]
\centering
\includegraphics[width=0.72\textwidth]{figs/check_cm20_dustem_ac.png}
\caption{Grain charging in the warm neutral medium: \texttt{ac} (circles) against DustEM's CM20 (lines). Top: mean
charge and its spread against radius. Bottom left: ratios of photoelectric current, gas current, absorbed power and
heating efficiency. Bottom right: charge distributions at four radii.}
\label{fig:charging}
\end{figure}

\paragraph{The PAH/grain seam.} At $\NC = 400$ the molecular tables (pahchem) meet the grain tables (\texttt{ac}). The
two should roughly agree there; they do not.
\begin{itemize}
\item Rerun on 12 November with pahchem's calibrated thresholds, PAHs heat more per C atom than \texttt{ac}: 0.9--2.7
  times in bin 1 and 1.9--3.6 times in bins 2--4. In CNM and dense gas, bins 2--4 are 40--75\,\% anions, against
  22--37\,\% for \texttt{ac}.
\item The October attribution of the gap: electron capture $\times$1.7--2.8, energy per electron $\times$1.3--1.9,
  absorption $\times$1.1--1.45, yield $\times$1.1.
\item The data side: PAHdb's DFT ionisation potentials sit 0.4--0.6\,eV below the measured ones. For cation
  recombination of small cations ($\NC$ 6--16), \texttt{ac} is within a factor 2 of laboratory data, while pahchem's
  default is 3--30 times too high.
\end{itemize}
Which side moves is an open decision (Sect.~\ref{sec:next}).

\begin{figure}[H]
\centering
\includegraphics[width=\textwidth]{figs/check_pah_grain_seam.png}
\caption{The PAH/grain seam. Left: charging thresholds against $\NC$ (dots: PAHdb with calibrated shifts; lines:
\texttt{ac}; dotted: DustEM). Middle and right: photoelectric heating $\Gamma$ and recombination cooling $\Lambda$ per C
atom for the four molecular bins, PAH tables against \texttt{ac}, in the WNM and a PDR.}
\label{fig:seam}
\end{figure}

%======================================================================================================================
\section{Hydrogen: tabulate or track?}

Montillaud+13 showed that hydrogenation is not always in equilibrium. C$_{66}$H$_{20}$ takes $10^2$\,yr to rehydrogenate
at $A_V = 0$ and $10^4$\,yr at $A_V = 1$ in NGC\,7023, while its charge relaxes in hours. Before adding scalars, we
therefore tested where an equilibrium table (E) is enough, where one tracked scalar (M) is needed, and where even one
scalar fails (F). The tests:
\begin{itemize}
\item T1, relaxation times;
\item T2, whether the charge can be eliminated;
\item T3, a one-scalar closure;
\item T4, parcels crossing PDRs at 0.3--3\,km\,s$^{-1}$, in and out of the cloud, and gas cycling between phases;
\item T5, mixing;
\item T6, a RAMSES-like frozen-rate step.
\end{itemize}

\subsection{Grains (laboratory laws)}
For the covalent grains the hydrogen balance has two terms (\texttt{gsd/hydrogen.py}). UV H loss is a laboratory yield
per absorbed photon (Alata+14, Mennella+01); re-hydrogenation by H atoms follows Mennella+02, with a plateau at
H/C\,$\approx$\,0.65. One scalar is then exact, and the step RAMSES would take, with rates frozen over a hydro step,
is exact to 0.05\,eV in $\Eg$. Re-hydrogenation turns out 10 times slower in the CNM than Mennella estimated, partly
because his balance used the flux $nv$ where $nv/4$ applies. The resulting treatment map:
\begin{itemize}
\item \textbf{E} in UV-lit gas changing on $\gtrsim 10^5$\,yr;
\item \textbf{M} in PDR transition layers, outflows from clouds, fast phase cycling and shielded gas;
\item \textbf{F} only where grains with different histories mix, for far-IR observables (a mixing bias of up to
  $-99$\,\% at 100\,\textmu m).
\end{itemize}
This is the case for rung M1 in the covalent bins (Figs.~\ref{fig:grainT1}, \ref{fig:grainpaths}).

\begin{figure}[H]
\centering
\includegraphics[width=0.80\textwidth]{figs/check_hydrogen_tracking_T1.png}
\caption{Grains, T1: relaxation time (top) and equilibrium band gap (bottom) of the 1--3\,nm bin against FUV and
$n$(H\,{\sc i}) at three temperatures. White lines are contours of constant relaxation time.}
\label{fig:grainT1}
\end{figure}

\begin{figure}[H]
\centering
\includegraphics[width=0.85\textwidth]{figs/check_hydrogen_tracking_paths.png}
\caption{Grains, T4: band gap of the 1--3\,nm bin along PDR paths, the true evolution at three speeds (blue) against the
equilibrium table (black), with the conditions below. Where the blue lines leave the black one, the grains remember
their history and one tracked scalar is needed.}
\label{fig:grainpaths}
\end{figure}

\subsection{PAHs: from ``track bin 1'' to ``tabulate everything''}
The molecular bins use pahchem's full model, with every PAHdb v4.00 molecule in each bin (712 in all). It took four
rounds to settle:
\begin{enumerate}
\item \textbf{Single parents} (C$_{24}$, C$_{54}$, C$_{96}$; 7 November). We reproduced Montillaud+13's Fig.~9 once the
  transition sits where theirs does. Near each molecule's own transition, with conditions changing on $\lesssim
  10^5$\,yr, neither E nor one scalar works, because the hydrogenation ladder is all-or-nothing.
\item \textbf{Full bins, charge-averaged chain} (8 November). Bin 1 fails E by 0.4--0.6 and is helped by one tracked
  $\langle\NH\rangle$; bins 2--4 are E. Proposal: $\MH$ in molecular bin 1 only.
\item \textbf{Full bins, joint $(Z, \NH)$ chain} (13 November), with super-hydrogenated states up to $N_{\rm H,0}+4$ and
  a charge for each $\NH$. Bin 1 becomes E through the PDRs. Bins 3--4 now fail in shielded PDR layers, through their
  super-hydrogenated states, and the +4 cap sets the answer there. The one-scalar closure fails too, because H addition
  alternates between slow and fast steps.
\item \textbf{Literature review and decision} (15 November). No code in the literature carries a hydrogenation state
  through moving gas, and what matters is UV-weighted and time-averaged. You decided: \emph{equilibrium tables, nothing
  tracked, judged on observables}.
\end{enumerate}

\paragraph{Observables.} \texttt{pahchem/hydrogen\_observables.py} gives each state $(Z, \NH)$ of each molecule a weight
per observable, $W[Z,\NH]$, so that $O = \sum W F$ holds for any state distribution $F$, steady or evolving. The
observables:
\begin{itemize}
\item carbon loss;
\item photoelectric heating and recombination cooling;
\item H$_2$ released and H taken up;
\item absorbed UV, together with the aromatic and aliphatic CH counts. Their UV-weighted ratio $r_{\rm ali}$ is the
  proxy for the 3.4/3.3\,\textmu m band ratio (Murga+22).
\end{itemize}
Carbon, heating, cooling and absorbed UV come from pahchem's own \texttt{joint\_state\_observables} (PR \#17), and
reproduce pahchem's steady-state functions to $10^{-9}$ or better.

\paragraph{Scoring.}
\begin{itemize}
\item \textbf{Paths:} 38, namely NGC\,7023, the Orion Bar and the Horsehead at 0.3, 1 and 3\,km\,s$^{-1}$ in and out;
  CNM$\leftrightarrow$WNM cycles; and $G_0\times10$ cycles at $n$(H)\,=\,10 with periods of $10^3$--$10^5$\,yr.
\item \textbf{Error measure:} on each path, E is compared with the time-dependent truth through the factor error
  $\max(E/T, T/E) - 1$.
\item \textbf{Worst cell and path budget:} the worst cell is taken over cells where the observable is significant. The
  path budget compares the path integrals, which is what an emitting or heating parcel feels.
\item \textbf{Carbon} is scored only where the bin loses its carbon within 1\,Gyr.
\end{itemize}

\begin{figure}[H]
\centering
\includegraphics[width=\textwidth]{figs/pah_paths.pdf}
\caption{PAH bins with the production chemistry (deep super-H chain, cap +8): the time-dependent truth (solid) against
the equilibrium tables (dashed) along four paths. Top row: the conditions. Then the photoelectric heating, carbon loss and
H$_2$ release per C atom, and the 3.4/3.3 proxy $r_{\rm ali}$. On the PDR paths the dashed lines lie on the solid ones.
The cycle with a $10^3$-yr period shows the lag of the larger bins.}
\label{fig:pahpaths}
\end{figure}

\begin{figure}[H]
\centering
\includegraphics[width=\textwidth]{figs/pah_scores.pdf}
\caption{Equilibrium tables against the truth: worst path-budget error per bin and observable, by kind of path (deep
chain, cap +8). Dashes: no carbon loss above the 1-Gyr floor.}
\label{fig:pahscores}
\end{figure}

\paragraph{Result} (Figs.~\ref{fig:pahpaths}, \ref{fig:pahscores}):
\begin{itemize}
\item On every PDR path, in every bin and for every observable, the tables are within 8\,\% of the truth in path budget.
\item Cycles with $10^4$--$10^5$\,yr periods stay within 0.5.
\item Only cycles with a $10^3$-yr period fail: carbon loss by a factor 2.2--2.4 in bins 1--2, and H$_2$ by up to 6.4.
  \begin{itemize}
  \item The H$_2$ error is real physics. Large PAHs lose their hydrogen more slowly than a 500-yr bright half-cycle,
    stay nearly intact, and shed H$_2$ from intact molecules.
  \item It sits in atomic gas, where PAH H$_2$ is about 1\,\% of grain formation and H$_2$ is photodissociated at once.
  \end{itemize}
\item Heating and cooling are within 11\,\% everywhere.
\end{itemize}

\subsection{The super-hydrogenation cap}
At the +4 cap, bins 3--4 sat at the top of the ladder in shielded gas. pahchem therefore built a deep chain
(\texttt{superh\_chemistry = "deep"}, PR \#18) with your choices of 15 November:
\begin{itemize}
\item the first four steps unchanged;
\item Boschman+15's alternating barriers beyond them, with Cazaux+16's barriers at the cation's +5, +11 and +17;
\item Eley--Rideal abstraction per extra H;
\item sequential H losses;
\item caps up to +24.
\end{itemize}
We reran all 712 molecules at +4 (the old chemistry) and with the deep chain at +8, +12 and +24 (Fig.~\ref{fig:cap}).
\begin{itemize}
\item \textbf{Heating, cooling and carbon loss do not depend on the cap or the chemistry:} $\le$\,2\,\% in path budget.
\item \textbf{Bins 1--2} have converged by +8.
\item \textbf{Bins 3--4 depend on the cap} for H$_2$, H uptake and $r_{\rm ali}$. Bin 4 has not converged even at +24.
  \begin{itemize}
  \item Below $G_0/n({\rm H}) \approx 10^{-2}$, past +4 the additions run on until the cap stops them.
  \item At +24, 76\,\% of bin 4 sits at the cap in the Orion Bar, and its aliphatic CH equal its aromatic CH on every
    path ($r_{\rm ali} \approx 1$).
  \end{itemize}
\end{itemize}
Whether that is allowed needs the band emission, not this proxy. The physical stop belongs to pahchem: barriers past
+24, the Eley--Rideal bracket ($\times2$, $\times7.5$), and the 3.4\,\textmu m band.

\begin{figure}[H]
\centering
\includegraphics[width=\textwidth]{figs/pah_cap.pdf}
\caption{The cap test along three paths (line styles) for the four bins (colours). Left: the path-averaged 3.4/3.3
proxy. Middle: the largest share sitting at the cap. Right: the path-averaged H$_2$ release relative to cap +24. A16 is
pahchem's default chemistry (Andrews+16) at +4; the others use the deep chain.}
\label{fig:cap}
\end{figure}

\begin{figure}[H]
\centering
\includegraphics[width=\textwidth]{figs/pah_T1_deep_k8.png}
\caption{PAH bins, T1 (deep chain, +8): the time the hydrogenation takes to settle (90\,\%) after a $\times10$ step in
$G_0$, with contours of the intact share. Dark regions settle in about an hour or less. The slow regions sit at low
density near each bin's transitions, which move to higher $G_0/n$ for larger molecules.}
\label{fig:pahT1}
\end{figure}

\subsection{The production tables}
You chose the deep chain at +8 for the first production build (16 November). It was built with
\texttt{pahchem/hydrogen\_tables.py} on megatron in 37\,min.
\begin{itemize}
\item \textbf{Grid:} the full pahchem-nd v1 grid. $\log G_0$ runs from $-1$ to 6 and $\log n$(H\,{\sc i}) from $-1$ to 6,
  both in 0.5 steps; $\log T$ from 1 to 4 and $\log x_e$ from $-5$ to $-1$; 7875 nodes in all.
\item \textbf{Files per bin:}
  \begin{itemize}
  \item \texttt{<bin>\_carbon.dat}, the carbon-loss table in the format RAMSES already reads;
  \item \texttt{<bin>.npz}, with every observable, the shares and the at-cap share.
  \end{itemize}
\item \textbf{Checks:} the files read back exactly, the shares sum to 1, and the build matches the pilot at shared nodes.
\item \textbf{Cap over the whole grid:} comparing +8 with +24 over the grid's corners, heating and cooling change by
  $\le$\,0.3\,\%. Carbon changes by $\le$\,0.13\,\% in bins 2--4. In bin 1 it changes by 33\,\% at one hot, dense node,
  through C$_{24}$H$_{14}$'s tiny dication share.
\item \textbf{Where the tables are provisional:} H$_2$ and $r_{\rm ali}$ are set by the cap wherever the gas is warm and
  shielded, in every bin. The at-cap share marks those nodes.
\item \textbf{Axis definition:} the density axis is the \emph{atomic} H density, and $x_e = n_e/n$(H\,{\sc i}).
\end{itemize}

\begin{table}[H]
\centering\small
\caption{Interpolation error of the production tables: log-linear interpolation against 286 nodes solved half a step
away.}\label{tab:interp}
\begin{tabular}{@{}l l@{}}
\toprule
observable & interpolation error \\
\midrule
heating, cooling & $\le$\,9\,\% \\
carbon loss, bin 1 & $\le$\,13\,\% \\
carbon loss, bins 2--3, $G_0 \le 10^{4.5}$ & 90th percentile $\le$\,0.2, worst $\times$1.7 \\
carbon loss, bins 2--3, $G_0 > 10^{4.5}$ & up to $\times$2.3 (bin 2) and $\times$24 (bin 3): 3.6\,dex per step \\
H$_2$ & up to $\times$9, where the addition barriers open between 10 and 32\,K \\
\bottomrule
\end{tabular}
\end{table}

\begin{figure}[H]
\centering
\includegraphics[width=0.92\textwidth]{figs/tables.pdf}
\caption{The production tables at $T = 100$\,K and $x_e = 10^{-4}$, one row per bin. Left: lifetime of the bin's carbon
(1\,Myr and 1\,Gyr contours). Middle: photoelectric heating per C atom. Right: mean hydrogenation $\langle\NH\rangle /
N_{\rm H,0}$: blue is dehydrogenated, red super-hydrogenated, and the black contour marks half the molecules at the
cap. Small PAHs are stripped and destroyed at high $G_0/n$; large ones survive everywhere but super-hydrogenate in
shielded gas.}
\label{fig:tables}
\end{figure}

\subsection{Why the grains are tracked and the PAHs tabulated}\label{sec:hwhy}
The literature does not say ``PAH hydrogenation is in equilibrium, small a-C(:H) is not'' (reading in
\texttt{lit/hydrogen\_grains\_vs\_pahs.md}):
\begin{itemize}
\item \textbf{PAH models.} Montillaud et al.\ (2013) find steady state reached ``from less than a year for small PAHs,
  up to $10^4$ years for large PAHs at $A_V = 1$'', because each multi-photon dehydrogenation step is a bottleneck, and
  conclude that ``large PAHs may never reach their steady state abundances''. Andrews et al.\ (2016) solve for steady
  state without testing it. Murga et al.\ (2022) run fixed environments for $10^5$\,yr without a steady-state
  comparison (our re-read of their chains: settled by $10^4$\,yr except in their molecular cloud).
\item \textbf{a-C(:H) models.} Jones et al.\ (2013) state that ``the dust processing time-scale is generally longer
  than the local dynamical time-scale and so the dust will therefore not be in equilibrium''; Jones et al.\ (2017):
  in equilibrium in diffuse and dense clouds, ``often out of equilibrium'' in PDRs and shocks. Murga et al.\ (2016,
  2019, 2023) follow the H content in time. Only the galaxy-scale one-zone model of Hirashita \& Murga (2020) takes it
  in equilibrium.
\end{itemize}
So our PAH decision does not claim the hydrogenation \emph{state} is in equilibrium: near each molecule's transition
it is not, as Montillaud et al.\ found and our single-parent tests reproduced. It claims that the PAH observables
RAMSES uses (heating, cooling, carbon loss) are within 8\,\% of the time-dependent truth on every PDR path.

\paragraph{The same metric for the grains.} The grain decision was made on the state (band gap). To compare like with
like, \texttt{diagnostics/gsd/score\_grain\_hydrogen\_obs.py} scores the grains as the PAHs were scored: equilibrium
against the truth, path budget of the per-gram opacity, on the same PDR paths and phase cycles
(Table~\ref{tab:grainobs}).

\begin{table}[H]
\centering\small
\caption{Grain hydrogen scored like the PAH tables: path-budget factor error of the equilibrium table against the
time-dependent truth, worst over the paths of each kind, for the 1--3\,nm bin (and the 20--53\,nm bin). ``Out'':
cloud material leaving the cloud (starting H-rich). 0.15\,\textmu m drives the photoelectric heating.}\label{tab:grainobs}
\begin{tabular}{@{}l l l l l@{}}
\toprule
paths & 0.15\,\textmu m & 0.55\,\textmu m & 3.4\,\textmu m & 100\,\textmu m \\
\midrule
Orion Bar, inward & 0.00 (0.00) & 0.02 (0.02) & 0.04 (0.07) & 0.13 (0.59) \\
Orion Bar, outward & 0.00 (0.03) & 0.11 (0.16) & 0.22 (0.59) & 1.2 (7.9) \\
Horsehead, inward & 0.01 (0.07) & 0.17 (0.59) & 0.60 (4.1) & 3.6 (104) \\
Horsehead, outward & 0.03 (0.19) & 1.6 (2.2) & 3.0 (8.7) & 250 (3100) \\
CNM$\leftrightarrow$WNM, $10^4$\,yr & 0.00 (0.01) & 0.02 (0.07) & 0.06 (0.30) & 0.35 (1.7) \\
\midrule
PAH bins, any PDR path (Fig.~\ref{fig:pahscores}) & \multicolumn{4}{l}{heating, cooling, carbon $\le$\,0.08} \\
\bottomrule
\end{tabular}
\end{table}

With the same metric, the grains' FUV absorption, and so their photoelectric heating, is as safe in equilibrium as the
PAHs' heating (within a few per cent on every path). What fails is what depends steeply on the band gap: the far-IR
opacity (factors of 2 to $10^3$ on the outward and Horsehead paths), the visible extinction and the 3.4\,\textmu m
region, mainly for material leaving clouds. The two decisions are consistent once they are stated per observable.

\paragraph{Why grains and PAHs differ physically.}
\begin{itemize}
\item \textbf{The observables differ in sensitivity.} PAH heating, cooling and carbon loss depend on the H state only
  for small, dehydrogenated molecules; elsewhere the charge sets them. The far-IR, visible and 3.4\,\textmu m opacity of
  a-C(:H) depend steeply on the band gap (0.1\,eV is a factor 2--3 in $\kappa_{100}$).
\item \textbf{A PAH's state is bimodal, a grain's is continuous.} Outside its narrow transition a PAH sits at one end of
  its ladder and gets there fast; the lag lives only near the transition. A grain's H/C is a bulk quantity anywhere
  between the UV floor and the H plateau.
\item \textbf{The slow step differs.} A PAH exchanges 10--20 edge H, and its slow step is the multi-photon H loss near
  the transition. A grain holds H/C of 0.05--0.65 over thousands of C, and its slow step is re-hydrogenation:
  $\sigma_f = 1.9\times10^{-18}$\,cm$^2$ (Mennella et al.\ 2002) is a net incorporation efficiency, $10^3$--$10^4$ below the
  geometric cross section, and Jones et al.\ (2014) find H incorporation slower than photolysis below $A_V\approx0.7$.
  Hence $10^5$--$10^6$\,yr relaxation in shielded, low-density gas.
\item \textbf{Memory.} Cloud material that leaves a cloud keeps its H-rich state for $10^4$--$10^5$\,yr.
\end{itemize}

\paragraph{Caveats.}
\begin{itemize}
\item The grain laws have no thermal (multi-photon) H loss, which would speed up the smallest grains; Jones \& Ysard
  (2025) treat only C--C bond loss. A rough estimate gives about 10\,\% H loss per photon above 10\,eV for sub-nm
  grains, which could move the 1--3\,nm bin towards tabulation: a sensitivity run is needed.
\item ``Outward'' paths start H-rich, an assumption (no C + H accretion physics yet).
\item Re-hydrogenation efficiencies are poorly constrained (Murga et al.\ 2023 call it ``undetermined'').
\item Elyajouri et al.\ (2024) fit $\Eg = 0.03$\,eV in the Orion Bar's atomic zone, below the laboratory floor.
\item The PAH bins and the smallest a-C bins may describe the same carriers (Jones 2012, paper III).
\item Our 2900\,yr relaxation at $G_0 = 1$ for 1--3\,nm and Jones \& Habart's (2015) $6.5\times10^4/G_0$\,yr are
  different quantities (a relaxation with re-hydrogenation against one-way photolysis, and different photon bands);
  still to be reconciled.
\end{itemize}

%======================================================================================================================
\section{Growth of the molecular bins: a bottom-up channel to decide}\label{sec:growth}

The plan of 30 October left molecular growth off by default ($f_{\rm acc} = 0$, as in PDR codes) and the bottom-up
source opt-in. A model of PAH evolution needs a bottom-up channel, so a choice has to be made now; the numbers can be
calibrated later.

\paragraph{Why the grain growth law cannot be used for molecules.} Applied to PAHs and nano-grains with sticking 1,
the classical law doubles a 50-C PAH's carbon in $10^3$--$3\times10^4$\,yr in the CNM, through Coulomb-focused
C$^+$ capture on anions, which contradicts the stable PAH sizes and gas-phase carbon of the diffuse ISM. Growth of
molecules must be reaction-limited: capture rate $\times$ the probability that the carbon stays $\times$ its survival
against UV.

\begin{table}[H]
\centering\small
\caption{What is known about the pieces of a bottom-up channel (references in the plan and
\texttt{reference\_data/carbon\_growth\_2026-10-30/}).}\label{tab:growthparams}
\begin{tabularx}{\textwidth}{@{}l L L@{}}
\toprule
piece & range & source \\
\midrule
C$^+$ capture $k_{\rm cap}$ & neutral $3$--$6\times10^{-9}$\,cm$^3$\,s$^{-1}$; anion $2.3\times10^{-7}\,T_2^{-1/2}$; cation $\approx0$ & Canosa+95 (anthracene), Langevin; Wolfire+08, Hollenbach+12 \\
fraction that adds carbon $f_{\rm acc}$ & 0 (PDR codes) to 0.4 (anthracene, the only measurement, $\NC = 14$) & Omont \& Bettinger 2021, 2025 \\
retention against UV $p_{\rm ret}$ & $\lesssim0.5$ below $\NC\approx40$--50 & Omont \& Bettinger 2025 \\
bound on net growth & $f_{\rm acc}p_{\rm ret}\lesssim0.15$ (slower than O sputtering) & Omont \& Bettinger 2021, App.~A.4 \\
C($^3$P) channel & $3$--$6\times10^{-10}$\,cm$^3$\,s$^{-1}$ (C + C$_6$H$_6$) $\times$ reactivity $r\in[0,1]$ & Hickson \& Loison 2026; Krasnokutski+17 \\
new molecules from gas & not derivable: TMC-1 small PAHs are 2.5--6\,dex above chemical models & McGuire+21; Wenzel+24/25; Byrne+24 \\
\bottomrule
\end{tabularx}
\end{table}

\paragraph{Why both a source and growth are needed.} Molecules made from the gas enter at $\NC\approx20$--24
(coronene is the largest PAH detected in TMC-1), below the $\approx$50-C photostability limit of the diffuse ISM; a
lasting population needs growth through bins 1--2 before the cloud disperses. Growth alone, without a source, only moves
carbon that shattering has already brought down. So the options are:

\begin{table}[H]
\centering\small
\caption{Options for the bottom-up channel. All add no scalars: growth is a reaction-limited $\dot a$ through the same
characteristic remap, the source an injection into bin 1 at its normal H/C.}\label{tab:growthopts}
\begin{tabularx}{\textwidth}{@{}l L L L@{}}
\toprule
option & what & calibrated on & expected weakness \\
\midrule
A & off (PDR codes): top-down only & -- & no TMC-1 PAHs; no translucent rise in $q_{\rm PAH}$ beyond shattering \\
B & growth only: C$^+$ and C capture with one efficiency $\epsilon = f_{\rm acc}p_{\rm ret}\le0.15$ & the $q_{\rm PAH}$ rise from $A_V$ 1 to 3 (Zhang+25) & no new molecules \\
C & source only: Form A, $\kappa\,n_{\rm H}n({\rm C})\,\Theta(A_V - A_c)$, or Form B, $\epsilon_{\rm ff}$ per free-fall time & TMC-1 & injected molecules stay small and are destroyed when exposed \\
D & source (Form A) + growth ($\epsilon$) & TMC-1 for $\kappa$; Zhang+25 for $\epsilon$ & two calibrated numbers \\
\bottomrule
\end{tabularx}
\end{table}

\paragraph{Recommendation: option D}, with option A kept as the ablation baseline and both numbers labelled
``calibrated'' in every output header.
\begin{itemize}
\item Form A rather than B as the default source: it lives where atomic carbon lives (the C\,{\sc i} layer that RTZ
  computes), scales with $n({\rm C})$ and so with metallicity, and vanishes in UV-lit gas without an extra cut. Form B
  stays for runs without atomic carbon.
\item One growth efficiency $\epsilon$ rather than the pair $f_{\rm acc}(\NC,Z)$, $p_{\rm ret}(\NC, G_0/n)$, until the
  data constrain more than one number; its upper limit is Omont \& Bettinger's 0.15.
\item Gates: the TMC-1 box (the bin-1 abundance within $\times10$); phase parity of $q_{\rm PAH}$ (1.0--1.1,
  Chastenet+19); at most $\times2$ from $A_V$ 1 to 3 (Zhang+25); $q_{\rm PAH}<1\,\%$ below 0.3\,Z$_\odot$ (Draine+07,
  Aniano+20, Whitcomb+24); the Sextans~A clumps (Tarantino+26); diffuse PAH sizes stable over Myr; gas-phase C at
  100--200\,ppm (Jenkins 2009; Sofia+11); and the budget limit (rate per gas-phase C atom $\times$ shielded fraction
  $\lesssim10^{-3}$ per Myr).
\end{itemize}
\textbf{To decide:} D (or another option) for the reference model; Form A or B; one $\epsilon$ or the two-parameter
form.

%======================================================================================================================
\section{How carbon moves between molecules, nano-grains, grains and the gas (23 November)}\label{sec:transfer}

Before A3c the shattering and coagulation modelling of carbon was reopened: what shattered fragments become, how
PAH-sized carbon moves up in size, and how erosion routes carbon between the species. The study has three
literature reviews (collisions; erosion; growth and identity) plus two on survival and release, a check of the
original design against them, and four quantitative checks (scripts \texttt{diagnostics/gsd/check\_carbon\_*.py} and
\texttt{check\_seam\_decomposition.py}; notes and outputs in \texttt{overview/lit/} and \texttt{overview/tables/}).

\paragraph{The original design stands where it was decided.} The 30 October design already routed nearly every
process, and the literature confirms it: fragments keep their parent's H/C (Murga et al.\ 2016's choice, the only
published one); molecules and nano-grains freeze out onto large grains (K\"ohler et al.\ 2015); growth of molecules is
reaction-limited; PAH carbon loss is a downward shift; sublimation uses graphite as the a-C proxy. The identity of carbon
below $\NC^{\max}$ is the hydrocarbon family along the hydrogen coordinate: a PAH unless bare (a-C-like) or fully
hydrogenated (a-C:H-like), the structure at equal H/C a labelled assumption. The literature cannot settle that
structure: size diagnostics fix the emitters' size (about 40--100\,C), not their planarity, and THEMIS and PAH-database
fits both reproduce the bands. Two records were found wanting: the October test behind ``3--4 bins over $\NC$
20--400 keep the 3.3\,\textmu m ratios within 2--7\,\%'' left no script or output (to redo with the extended
$\NC^{\max}$ test), and $\NC^{\max} = 400$ loses its main anchor (evaporating VSGs as clusters of $\ge400$\,C) with the
clusters; the band physics still supports 150--400\,C.

\paragraph{What the checks show.}
\begin{itemize}
\item \textbf{The seam is first-order for the gas.} For THEMIS J13 carbon (43\,\% of it below 400\,C) that carbon carries
  70--80\,\% of the grain photoelectric heating and recombination cooling; charging it with the PAH tables instead of
  \texttt{ac} multiplies the total heating by 1.4 (CNM, dense) to 2.5 (WNM, PDR), and still by 1.2--2.0 if only
  10--20\,\% of the carbon is below 400\,C (\texttt{check\_carbon\_heating\_share.py}).
\item \textbf{Timescales} (\texttt{check\_carbon\_transfer\_timescales.py}): shattering refills the small bins in about
  $10^7$\,yr in warm gas; bin 1 (20--42\,C) photolyses in $2\times10^6$\,yr (diffuse) and 300\,yr (PDR), bins 3--4 never;
  a-C hydrogen relaxes in $3\times10^3$\,yr (diffuse) and 3\,yr (PDR) with the laboratory laws, so the equilibrium band
  gap is consistent where there is UV; freeze-out onto grains takes $3\times10^7$\,yr (CNM) and $3$--$5\times10^5$\,yr
  (dense cloud; Seok et al.\ 2014: 0.3--0.5\,Myr).
\item \textbf{Molecules must not collide with molecules under grain physics.} In the code as it was, PAH--PAH sticking
  (cold gas) and even PAH--PAH ``shattering'' (warm gas: thermal speeds of $\sim$3\,km\,s$^{-1}$ above graphite's
  1.2\,km\,s$^{-1}$) dominate the molecular bins 10--25 times over freeze-out.
\item \textbf{One-zone runs} (\texttt{check\_carbon\_transfer\_options.py}, Table~\ref{tab:transfer}): whether molecules
  stick to nano-grains decides how much PAH-sized carbon survives in cold gas.
\end{itemize}

\begin{table}[H]
\centering\small
\caption{One-zone runs of the carbon grid from THEMIS J13 (collisions with graphite strength and HY09 sticking, neutral
grains; molecular carbon loss from the pahchem tables): carbon shares at the end. O0: the code as it was; O1: no
molecule--molecule outcome, molecules only stick to larger particles; O2: as O1, molecules stick only to grains
$\ge20$\,nm. Last column: total photoelectric heating with the molecular bins on PAH tables over \texttt{ac}.}\label{tab:transfer}
\begin{tabular}{@{}l l r r r r r@{}}
\toprule
cell & rule & molecular & nano 1--20\,nm & grains $\ge20$\,nm & gas & heating PAH/ac \\
\midrule
start & & 0.463 & 0.272 & 0.266 & 0 & \\
WNM, 30\,Myr & O0 & 0.283 & 0.165 & 0.002 & 0.550 & 2.77 \\
 & O1 / O2 & 0.348 & 0.309 & 0.014 & 0.329 & 2.62 \\
CNM, 30\,Myr & O0 & 0.038 & 0.542 & 0.256 & 0.164 & 1.10 \\
 & O1 & 0.087 & 0.444 & 0.259 & 0.210 & 1.25 \\
 & O2 & 0.246 & 0.273 & 0.261 & 0.220 & 1.57 \\
dense, 1\,Myr & O0 & 0.004 & 0.670 & 0.326 & 0 & 1.01 \\
 & O1 & 0.028 & 0.649 & 0.323 & 0 & 1.07 \\
 & O2 & 0.398 & 0.251 & 0.351 & 0 & 1.43 \\
\bottomrule
\end{tabular}
\end{table}

\paragraph{Observations of PAH survival.} PAHs survive at molecular-cloud scale: $q_{\rm PAH}$ in molecular gas equals
the diffuse value to $\pm30$\,\% in the LMC and SMC (Chastenet et al.\ 2019), JWST finds the PAH abundance insensitive to
the molecular fraction (Chastenet et al.\ 2023), and translucent clouds show no drop from $A_V$ 1 to 3 (Zhang et al.\
2025). Dense, dark filaments and cores ($n_{\rm H}\gtrsim$ a few $10^3$, $A_V\gtrsim2$--3) lose small-grain emission by a
factor of a few (Stepnik et al.\ 2003; Flagey et al.\ 2009; Ysard et al.\ 2013), partly from the weaker heating field.
Sticking to everything conflicts with the first (PAH bins 2--4 lose 72\,\% in 30\,Myr of CNM); sticking only to big
grains (21\,\%) passes it but may lose too little in cores.

\paragraph{Release from nano-grains.} A PAH binds to a carbon surface at about 0.05\,eV per C (Zacharia et al.\ 2004).
One UV photon heats a 1\,nm host to 400--640\,K but a 3\,nm host only to 100--140\,K, so thermal release is a
single-photon effect for small PAHs on the smallest hosts (Montillaud \& Joblin 2014; Lange et al.\ 2023). No other
channel was found that releases a PAH from a host of 3\,nm or more (cosmic-ray heating, reactive desorption, collisions,
host photo-fragmentation, shocks: all negligible or $\gtrsim10^8$\,yr) except non-thermal photodesorption, unmeasured for
PAHs on carbon or ice: a yield of only $\sim10^{-7}$ per absorbed photon balances freeze-out in the CNM, while a dense
cloud would need $\sim10^{-3}$. Such a yield alone gives the observed pattern.

\paragraph{The seam is electron sticking.} Redone with today's pahchem (four-state chains reproducing both full solves),
the heating gap decomposes as: photoemission rate $\times$0.95--1.4, energy per photoelectron $\times$0.7--1.3, ion
channels $\times$1.0--1.2, and electron capture $\times$1.3--2.7 in bins 2--4. Both sides use the same Coulomb focusing;
the difference is the sticking on cations: pahchem's laboratory fit reaches 1 by $\sim$35\,C, WD01's grain law gives
0.2--0.3. No recombination rate has been measured beyond 16\,C; WD01's 0.5 ceiling and 1\,nm escape length are
assumptions; C$_{60}$'s neutral attachment ($\sim$1, Viggiano et al.\ 2010) is the only point above. The plausible range
is 0.3--1 for 35--400\,C, and the gap lies inside it.

\paragraph{Decided (23 November).}
\begin{itemize}
\item No PAH clusters (rung M2's cluster part and P4 removed); the legacy CALIMA PAH--PAH and PAH--grain methods are not
  reused (they stay untouched for runs without the GSD).
\item Molecule + molecule: no outcome. Molecules never shatter or crater; they only stick to larger particles.
\item Molecule + grain $\ge20$\,nm: unit sticking (freeze-out). Molecule + nano-grain: $S_{\rm eff} = 1/[1 + R_{\rm abs}
  (P_{\rm th} + Y_{\rm pd})\,t_{\rm bury}]$, with $R_{\rm abs}$ the adsorbed molecule's photon absorption rate, $P_{\rm th}$
  the single-photon thermal release (derived), $Y_{\rm pd}$ the non-thermal photodesorption yield (an input labelled
  ``calibrated'', default 0, calibrated in A5 against the parity and dense-core gates, expected $10^{-7}$--$10^{-6}$) and
  $t_{\rm bury}$ the host's burial time.
\item The seam's electron sticking: one labelled law, pahchem's laboratory fit for the molecules, \texttt{ac} taking the
  same value at $\NC = 400$ and tapering to WD01 by $\sim$3\,nm; A5 runs the brackets (WD01 everywhere; the molecular
  value everywhere).
\end{itemize}

\paragraph{Open.}
\begin{itemize}
\item \textbf{Hot-gas erosion of the smallest carbon.} Classical sputtering fails below $\NC\approx1000$ (Bocchio et al.\
  2012). pahchem has modern molecular rates (thermal electron and ion passages, C and H knock-outs; \texttt{population.hot\_passages},
  \texttt{hot\_ion\_passages}), but the production tables stop at $10^4$\,K. Proposal: extend the tables' temperature axis
  and join pahchem's rates to the classical yields at $\NC\approx1000$, where Bocchio et al.\ find the two agree.
\item \textbf{Coulomb explosion: fragments or ion loss.} Two limits compete (Waxman \& Draine 2000; the code's
  \texttt{explosion\_rate}): mechanical failure when $(U/a)^2/4\pi$ reaches the tensile strength $S_{\max}$ tears the grain
  into fragments, while ion field emission at $\sim3\times10^8$\,V\,cm$^{-1}$ sheds ions one by one. The crossover is at
  $S_{\max}\approx8\times10^{10}$\,erg\,cm$^{-3}$: flawed or porous grains ($10^9$--$10^{10}$) explode, defect-free grains
  near the ideal strength ($\sim10^{11}$) shed ions, as the other dissociation channels do. The routing follows from the
  assumed strength, known only for bulk material; the fragments' size law ($a^{-3.3}$) is an assumption.
\item The a-C:H end of the hydrocarbon family (P1) and the bare end (pahchem's bare states or a-C physics).
\item Labelled omissions: vaporisation above $\sim$19\,km\,s$^{-1}$, fullerenes, cosmic-ray processing.
\end{itemize}

%======================================================================================================================
\section{What RAMSES will need}

\paragraph{Memory.} Measured on today's RAMSES-CALIMA: every table is real*8 and read in full by every MPI rank, at
11--13\,MiB per bin. Collisional cooling and sputtering dominate, then the WDB06rt charging kernels. In today's formats
the GSD scheme would need about 190\,MiB per rank: 8 carbon and 8 silicate grain bins, 4 molecular bins, and the
\texttt{acx} nodes. The plan brings that to about 20\,MiB per rank, under the 46\,MB target, with five principles:
\begin{enumerate}
\item export only expensive integrals needed per cell: photoemission yields, sputtering yields, multiphoton cascades
  (the PAH tables above), Mie and Jones optics;
\item compute in RAMSES whatever is analytic or a cheap combination of tables already in memory: accretion,
  coagulation, the charge solve, the grain H step;
\item never load post-processing quantities, such as band emission and heat capacity; those go to the exporter for
  SKIRT;
\item use axes the physics is smooth in, with real*4 storage. Uniform coarsening of the cooling and sputtering tables
  fails, so they are re-gridded on $(\log T, \phi/kT)$ first;
\item keep one copy per node, through MPI-3 shared memory; this is the largest single lever.
\end{enumerate}

\paragraph{Phase B in RAMSES} (feature branch off \texttt{staging}; to be created only when you say so):
\begin{itemize}
\item \textbf{B1, state:} the carbon $\MC$ block, then the $N$ block, then the $\MH$ block. \texttt{NGSD=0} keeps today's
  behaviour byte for byte.
\item \textbf{B2, sources per cell in \texttt{dust\_solver.f90}:}
  \begin{itemize}
  \item first the exact H step: $\MH \to M_{\rm eq} + (\MH - M_{\rm eq})\,e^{-(k_{\rm UV}+k_{\rm H})\Delta t}$, with
    $k_{\rm UV}$ from the RT group photons and $k_{\rm H}$ from $n$(H\,{\sc i}) and $T$;
  \item then the carbon-moving processes, each carrying H/C;
  \item the grains' rates on the gas: one charge solve per bin at its area node (Sect.~\ref{sec:inbin}), the
    surface processes from the closure's second moment;
  \item one overflow bin per composition above $a_{\max}$, computed only above $10^{-6}$ of the cell's dust mass,
    with its share logged and a warning above 1\,\% (Sect.~\ref{sec:mck}).
  \end{itemize}
\item \textbf{Transport:}
  \begin{itemize}
  \item a ratio pass at the CMA hook, with $N$ outside the CMA groups;
  \item $\MH$ as a CMA child of $\MC$, its slopes on $x$ = H/C;
  \item donor-ratio TVA fluxes for $N$ and $\MH$;
  \item an HLLC, exact or acoustic Riemann solver;
  \item recomputing the TVA drift (B4) instead of storing it.
  \end{itemize}
\item \textbf{B6/B7, tables per GSD bin:} optics, band emission and charging, molecular- or grain-type from the header.
  The RT optics come as five pseudo-bins per bin: ordinary dust bins for the SED table, the star polling and the IR
  means, blended linearly in the cell's $s$ once per cell-call. css is dropped (Sect.~\ref{sec:rtgroups}).
\item \textbf{Targets:}
  \begin{itemize}
  \item a CALIMA step at most 1.5 times the legacy one;
  \item Godunov at most 35\,\% slower;
  \item tables $\le$\,46\,MB per rank;
  \item the nightly suite byte-identical with \texttt{NGSD=0}.
  \end{itemize}
\end{itemize}

%======================================================================================================================
\section{Where we go from here}\label{sec:next}

\subsection{Immediate: consolidate}
\begin{enumerate}
\item \textbf{Done (19--22 November):} the GSD branch is merged (PR \#20), with the charge solve ported from RAMSES
  (PR \#19); the checks were rerun on \texttt{main} after the collision solver's fix (PR \#22), see ``Rerun on 22
  November'' in the summary. After that rerun EA misses 5\,\% at 8 bins only in the shattering case (7.1\,\%; under
  10\,\% everywhere); 12 bins would bring every case under 5\,\%. Decided (22 November): keep 8 bins.
\item \textbf{Re-measure the collision variants} (kernel averaging, Gauss nodes) through \texttt{GSDMethod} on the fixed
  solver: the variant bench of Sect.~\ref{sec:variants} predates that API and no longer runs.
\item \textbf{Merge pahchem's PRs \#17 and \#18,} then work from \texttt{main}. Rerun the checks that read
  pahchem (seam, bins), as every pahchem merge so far has moved them. The production tables were built from the
  unmerged code and need a rebuild only if pahchem changes before the merge.
\end{enumerate}

\subsection{Decisions waiting for you}
\begin{itemize}
\item \textbf{The bottom-up channel} for the molecular bins (Sect.~\ref{sec:growth}): recommended, a Form-A source
  calibrated on TMC-1 plus one reaction-limited growth efficiency calibrated on the translucent $q_{\rm PAH}$ rise.
\item \textbf{Drift and size-dependent sinks per moment} (Sect.~\ref{sec:diallo}): recommended, $N$ and $M$ with their
  own power-law velocities and rates, as RAMSES-MCR does for cosmic rays.
\item \textbf{Kernel averaging} in the collisions (Sect.~\ref{sec:variants}): available as an option; to be tested in
  the full collision model before deciding on a default.
\item \emph{Decided 23 November} (Sect.~\ref{sec:transfer}): the carbon mass transfer. No PAH clusters; molecules
  do not interact with molecules; they freeze out onto grains $\ge20$\,nm everywhere and onto nano-grains with a
  photo-release term whose non-thermal yield is a labelled, calibrated input (default 0); the seam's electron
  sticking is one labelled law (molecular below $\NC = 400$, tapering to WD01 by $\sim$3\,nm), bracketed in A5.
\item \textbf{Still open from the transfer study:} PAH erosion in hot gas with pahchem's impact rates (the tables stop
  at $10^4$\,K) joined to classical sputtering at $\NC\approx1000$; the routing of Coulomb explosion (fragments or ion
  loss, set by the assumed tensile strength); the a-C:H end of the hydrocarbon family (P1) and the bare end; redoing
  the band test.
\item \textbf{The super-H physical stop} for large PAHs, a pahchem question: barriers past +24, the Eley--Rideal bracket,
  a check against the 3.4\,\textmu m band. Until then, the H$_2$ and $r_{\rm ali}$ tables of bins 3--4 stay provisional.
\item \textbf{A finer $G_0$ axis above $10^4$} for the carbon-loss tables of bins 2--3, if the runs reach fields near
  massive stars. The format allows it, and the rebuild takes under an hour.
\item \emph{Decided 22 November:} $a_{\rm H,max}$. The four covalent bins below 53\,nm use the equilibrium band
  gap, nothing tracked by default; $\MH$ is compile-time for 0--4 bins (default 0) and A5 decides, bin by bin. Heating
  and charging are safe in equilibrium (Table~\ref{tab:grainobs}); what fails (visible, 3.4\,\textmu m, far-IR for
  material leaving clouds) is mostly post-processing; a tracked bin costs 16\,B per cell and one transported scalar,
  while the band-gap tables are needed either way. 8 bins per composition kept after the collision fix (same day).
\item \emph{Decided 22 November:} the carbon material. THEMIS a-C(:H) at 1.6\,g\,cm$^{-3}$ is the default and
  graphite the control (Sect.~\ref{sec:materials}); a material is a table set, so the choice stays reversible.
\item \textbf{The default $\NC^{\max}$} (400) after the extended band test; whether to use a $\xi(a)$ transition band
  for the DL07 comparison.
\item \textbf{Smaller choices:} H$_2$ release from cold grains (diffusion-limited or immediate); whether sputtering
  removes H first; associative detachment on intact PAH anions, which decides the dense-gas PAH charge.
\item \textbf{Which observational datasets to fetch first} for T3/T4: extinction (Gordon+23), $\tau_{3.4}/A_V$, the JWST
  PDR nano-grain results, Pilleri+15 maps, and others.
\end{itemize}

\subsection{The build, in order}
\begin{enumerate}
\item \textbf{Finish A3b:}
  \begin{itemize}
  \item the grain rates over the in-bin profile (two nodes per bin, or the chosen method), with
    \texttt{check\_inbin\_rates} as the test;
  \item charge and photoelectric heating per GSD bin from the PAH group tables, conditional on the H state;
  \item molecular band emission per bin, for post-processing;
  \item the re-gridded, real*4 cooling and sputtering tables.
  \end{itemize}
\item \textbf{A3c, the \texttt{ramses} mode of \texttt{gsd}:} the per-cell operator RAMSES will mirror. It reuses the
  remap, collision, sputtering and closure modules (EA default), and adds the $\MH$ block with the exact step, and H carried by every
  remap, a carried overflow, and moment-specific rates for size-dependent sinks. Its conservation tests (carbon and
  hydrogen to round-off, bounds on H/C) come first; a third collision node or bin edges at kernel thresholds can be
  tested there (the coagulation test of Sect.~\ref{sec:mck}). Also: a thermal H-loss sensitivity run for the 1--3\,nm
  bin (Sect.~\ref{sec:hwhy}).
\item \textbf{A4, advection prototype:} $N$ following $\MC$ through the ratio pass, $\MH$ as a child of $\MC$, and
  the drift with moment-specific velocities (Fig.~\ref{fig:drift} is its first test).
\item \textbf{A5, accept or reject.} T2: our tables against DustEM's SED and extinction, and the graphite control
  against DL07. T3: the diffuse ISM per H, the bump, $\tau_{3.4}/A_V$, the growth gates. T4: one-zone and slab runs for
  NGC\,7023, the Orion Bar and the Horsehead. The convergence of the thermal and chemical rates with the bin number
  (Sect.~\ref{sec:rateconv}). Then the ablations that decide the rungs, and the growth gates once the bottom-up
  channel is chosen.
\item \textbf{A6, the exporter} \texttt{models/export\_gsd\_tables.py}, with provenance headers; memory estimated with one
  worker first.
\item \textbf{Phase B} in RAMSES (above).
\end{enumerate}

%======================================================================================================================
\appendix
\section{Where things live}
\small
\begin{tabularx}{\textwidth}{@{}l L@{}}
\toprule
\multicolumn{2}{@{}l}{\textbf{code} (\texttt{src/pycalima/})}\\
\texttt{gsd/} & \texttt{closures} (the methods: EA default, EAF, PL2M, affine, exp), \texttt{optics} (incl.\ node tables), \texttt{pl2m}, \texttt{remap}, \texttt{collisions}, \texttt{sputtering}, \texttt{grid} (graded),
  \texttt{materials}, \texttt{carbon} (\texttt{CarbonLayout}), \texttt{optics}, \texttt{charging}, \texttt{hydrogen} \\
\texttt{models/dust\_radiation/} & \texttt{themis\_optics}, \texttt{jones2012\_optics}, \texttt{carbon\_heat\_capacity};
  \texttt{models/tools/coated\_mie} \\
\texttt{pahchem/} & \texttt{hydrogen\_tracking}, \texttt{hydrogen\_observables}, \texttt{hydrogen\_tables} (new, uncommitted) \\
\midrule
\multicolumn{2}{@{}l}{\textbf{checks} (\texttt{diagnostics/gsd/})}\\
& \texttt{check\_bin\_convergence}, \texttt{check\_optics\_tables}, \texttt{check\_carbon\_carriers},
  \texttt{check\_hydrogenated\_tables}, \texttt{check\_pah\_grain\_seam}, \texttt{check\_hydrogen\_tracking},
  \texttt{check\_pah\_hydrogen\_tracking}, \texttt{check\_pah\_bins\_hydrogen}, \texttt{compare\_pah\_bins\_caps},
  \texttt{analyse\_pah\_bins\_tracking}; \texttt{diagnostics/dust\_charge/check\_cm20\_dustem}; new on 17 November:
  \texttt{check\_inbin\_rates}, \texttt{check\_rates\_convergence}, \texttt{check\_mckinnon18\_tests},
  \texttt{check\_inbin\_drift}, \texttt{score\_grain\_hydrogen\_obs}; the variant study: \texttt{pl2m\_variants},
  \texttt{bench\_pl2m\_variants}, \texttt{bench\_pl2m\_figures} (results and \texttt{REPORT.md} in
  \texttt{gsd\_overview/variants/}); node tables and RT pseudo-bins: \texttt{check\_node\_tables},
  \texttt{check\_rt\_group\_optics}, \texttt{bench\_rt\_group\_optics} \\
\midrule
\multicolumn{2}{@{}l}{\textbf{documentation} (\texttt{docs/physics/})}\\
& \texttt{carbon-optics}, \texttt{charging-cm20} (incl.\ the seam), \texttt{grain-hydrogen}, \texttt{pah-hydrogen} \\
\midrule
\multicolumn{2}{@{}l}{\textbf{results} (\texttt{results/}, git-ignored)}\\
& \texttt{binconv\_v4/} (A0, A2), \texttt{optics/}, \texttt{charging/}, \texttt{carbon\_plan/} (hydrogen, seam, PAH bins
  and caps, production tables in \texttt{hydrogen\_tables/deep\_k8/}), \texttt{gsd\_overview/} (this report and
  \texttt{make\_figures.py}); literature notes for this revision in \texttt{gsd\_overview/lit/}; the papers in
  \texttt{reference\_data/gsd\_solvers\_2026-11-17/} \\
\midrule
\multicolumn{2}{@{}l}{\textbf{megatron}}\\
& \texttt{\textasciitilde/pycalima\_deep} (worktree plus pahchem PRs \#17, \#18), runs in
  \texttt{\textasciitilde/pycalima\_runs/}; launch with \texttt{setsid nohup nice -n 19} \\
\bottomrule
\end{tabularx}

\section{Glossary}
\begin{tabularx}{\textwidth}{@{}l L@{}}
PL2M & two-moment power-law closure: $N$ and $M$ per bin, a power law $dN/dm\propto m^\beta$ inside it (comparison method) \\
$s$ & in-bin state, $(\bar m/m_L - 1)/(\eta - 1)$ with $\bar m = M/N$: 0 with all grains at the lower edge, 1 at the upper \\
two nodes & the two quadrature points per bin (number halves at their mean masses) that reproduce $N$ and $M$ exactly \\
PC, PL & McKinnon et al.'s piecewise-constant and piecewise-linear (in $a$) bins \\
MPRK22 & modified Patankar Runge--Kutta, second order: positive and conservative for any step \\
EA, EAF & edge-aware closures: a bin next to an (almost) empty one is a front, cut off inside the bin; EA (the default) continues the neighbour's slope, EAF fixes it at $-3.5$ \\
node tables & per-gram quantities at a few radii per bin, integrated over the cell's reconstruction (any closure) \\
pseudo-bins & a bin's PL2M per-gram optics at 5 fixed $s$, each an ordinary RAMSES dust bin, blended linearly in the cell's $s$ \\
area node & the one radius per bin where RAMSES solves the charge, $a_* = \int a^3dN/\int a^2dN$, weighted to the bin's exact grain area \\
overflow bin & the open-ended bin above $a_{\max}$ ($N$, $M$; grains at its mean size), computed only when it holds grains \\
KA & kernel averaging over sub-slices of each node pair in the collision operator \\
$\eta$ & a bin's mass ratio $m_R/m_L$ \\
E / M / F & equilibrium table enough / one tracked scalar needed / one scalar not enough \\
joint chain & pahchem's network over charge $Z$ and hydrogen number $\NH$ together \\
super-H, cap & states with more H than the parent PAH ($\NH > N_{\rm H,0}$); the cap is the highest $N_{\rm H,0}+K$ allowed \\
$r_{\rm ali}$ & UV-weighted aliphatic over aromatic CH: the 3.4/3.3\,\textmu m proxy \\
path budget & factor error between the path integrals of E and the truth \\
TVA & terminal-velocity approximation: the dust drift relative to the gas, one stored drift per grain species \\
CMA & consistent multi-fluid advection, which keeps passive-scalar groups summing correctly \\
\texttt{ac}, \texttt{ach}, \texttt{cm20}, \texttt{acx} & charging materials: a-C (default), a-C:H, DustEM's CM20, intermediate H/C nodes \\
\end{tabularx}

\section{References}
{\footnotesize
\setlength{\parindent}{0pt}
\newcommand{\rf}[1]{\par\hangindent=1.5em #1}
Identifiers checked for this revision (arXiv pages or full texts; notes in \texttt{gsd\_overview/lit/}). Papers cited
only through the plan (growth chemistry, laboratory hydrogenation) are in \texttt{reference\_data/} with the plan's
notes and are not repeated here.

\medskip\textbf{Size-distribution solvers}
\rf{Aoyama S., Hirashita H., Nagamine K. 2020, MNRAS 491, 3844 (arXiv:1906.01917)}
\rf{Choban C. R., Kere\v{s} D., Hopkins P. F., Sandstrom K. M. 2022, MNRAS 514, 4506 (arXiv:2201.12369)}
\rf{Choban C. R., Salim S., Kere\v{s} D., Roman-Duval J., Sandstrom K. M. 2026, MNRAS 549, stag1020 (arXiv:2603.08504)}
\rf{Diallo N.-S., Dubois Y., Marcowith A., Rosdahl J., Commer\c{c}on B. 2026, RAMSES-MCR (arXiv:2602.21147)}
\rf{Dubois Y. et al.\ 2024, A\&A (arXiv:2402.18515)}
\rf{Hirashita H. 2015, MNRAS 447, 2937 (arXiv:1412.3866)}
\rf{Hirashita H., Aoyama S. 2019, MNRAS 482, 2555 (arXiv:1810.07962)}
\rf{Hirashita H., Yan H. 2009, MNRAS 394, 1061}
\rf{Hopkins P. F. 2023, the spectral cosmic-ray correction factors as used by Diallo et al.\ (2026, Sect.~2.2.5)}
\rf{Kopecz S., Meister A. 2018, BIT Numer.\ Math.\ 58, 691 (MPRK22)}
\rf{Lebreuilly U., Commer\c{c}on B., Laibe G. 2019, A\&A 626, A96 (arXiv:1905.01948)}
\rf{Lebreuilly U. et al.\ 2023 (arXiv:2211.01891)}
\rf{Li Q., Narayanan D., Torrey P., Dav\'e R., Vogelsberger M. 2021 (arXiv:2012.03978)}
\rf{Lombart M., Laibe G. 2021, MNRAS 501, 4298 (arXiv:2011.12298)}
\rf{Matsumoto K. et al.\ 2026, A\&A 705, A75 (arXiv:2508.21157)}
\rf{McKinnon R., Vogelsberger M., Torrey P., Marinacci F., Kannan R. 2018, MNRAS 478, 2851 (arXiv:1805.04521)}
\rf{Narayanan D. et al.\ 2023, ApJ 951, 100 (arXiv:2301.07136); Narayanan D. et al.\ 2026 (arXiv:2606.20809)}
\rf{Rodr\'iguez Montero F., Dubois Y., Katz H., Slyz A., Devriendt J. 2026, CALIMA (arXiv:2602.21790)}
\rf{Romano L. E. C., Nagamine K., Hirashita H. 2022, MNRAS 514, 1441 (arXiv:2202.05521)}
\rf{Sportisse B. 2000, J.\ Comput.\ Phys.\ 161, 140}

\medskip\textbf{Grain growth in dense gas (the overflow bin)}
\rf{Pagani L., Steinacker J., Bacmann A., Stutz A., Henning T. 2010, Science 329, 1622 (doi:10.1126/science.1193211; arXiv:1110.4180)}
\rf{Steinacker J., Pagani L., Bacmann A., Guieu S. 2010, A\&A 511, A9 (arXiv:0912.0145)}

\medskip\textbf{Carbon materials and the molecular limit}
\rf{Bern\'e O. et al.\ 2007, A\&A 469, 575 (arXiv:astro-ph/0703072)}
\rf{Chiar J. E. et al.\ 2006 (arXiv:astro-ph/0607245)}
\rf{Croiset B. A., Candian A., Bern\'e O., Tielens A. G. G. M. 2016, A\&A 590, A26 (arXiv:1603.02577)}
\rf{Draine B. T. 2003, ARA\&A 41, 241 (arXiv:astro-ph/0304489)}
\rf{Draine B. T., Li A. 2007, ApJ 657, 810 (arXiv:astro-ph/0608003)}
\rf{Draine B. T. et al.\ 2021, ApJ 917, 3 (arXiv:2011.07046)}
\rf{Fanciullo L. et al.\ 2015, A\&A 580, A136 (arXiv:1506.07011)}
\rf{Hensley B. S., Draine B. T. 2023, ApJ 948, 55 (arXiv:2208.12365)}
\rf{J\"ager C., Huisken F., Mutschke H., Llamas Jansa I., Henning Th.\ 2009, ApJ 696, 706 (arXiv:0903.0775)}
\rf{Jones A. P. 2012, A\&A 540, A1; 540, A2; 542, A98 (arXiv:1511.01673, 1511.01682, 1511.01684)}
\rf{Jones A. P. et al.\ 2013, A\&A 558, A62 (arXiv:1411.6293); Jones A. P. et al.\ 2017, A\&A 602, A46 (arXiv:1703.00775)}
\rf{Pilleri P. et al.\ 2015, A\&A 577, A16 (arXiv:1502.04941)}
\rf{Planck Collaboration Int.\ XXIX 2016, A\&A 586, A132 (arXiv:1409.2495)}
\rf{Rapacioli M. et al.\ 2005, A\&A 429, 193; Rapacioli M., Calvo F., Joblin C. et al.\ 2006, A\&A 460, 519 (arXiv:astro-ph/0609186)}
\rf{Ricca A., Bauschlicher C. W., Boersma C., Tielens A. G. G. M., Allamandola L. J. 2012, ApJ 754, 75}

\medskip\textbf{Hydrogenation}
\rf{Andrews H. et al.\ 2016, A\&A 595, A23}
\rf{Elyajouri M. et al.\ 2024 (arXiv:2401.01221)}
\rf{Hirashita H., Murga M. S. 2020 (arXiv:2001.01968)}
\rf{Jones A. P., Habart E. 2015 (arXiv:1507.08534); Jones A. P. et al.\ 2014, Faraday Discuss.\ (arXiv:1411.5877)}
\rf{Jones A. P., Ysard N. 2025 (arXiv:2511.04750)}
\rf{Montillaud J., Joblin C., Toublanc D. 2013, A\&A 552, A15 (arXiv:1301.6507)}
\rf{Murga M. S. et al.\ 2016 (arXiv:1612.00419); 2019 (arXiv:1906.11308); 2022, MNRAS 509, 800 (arXiv:2111.05583); 2023 (arXiv:2212.04993)}
}

\end{document}
